%TODO 等的使用
%%TODO 内容问题
%%FIXME 格式问题
%%HACK 粗鄙技巧
%%CODE 代码改进

%TODO——注意事项
%FIXME：20160422 en dash 与连接号 hyphen （人名中的使用）
%FIXME：20160621 adobe-garamond 的 semibold 与 bold 使用时需注意。实现方法是将字体文件重命名
%FIXME：20160625 由于改用 TeX Live 2016，上一条可能失效
%FIXME：20160625 英文名词改用半角括号
%FIXME：20160625 Hamilton 函数使用花体 H（\scH）
%CODE：20160625 每行 80 字符限制，中英文、数字之间加空格
%% 理想气体、可逆、绝热等公式特殊标注
%% 绝热：Adiabatic 等压：Isobaric 等容：Isochoric 等温：Isothermal 等熵：Isentropic 等焓：Isenthalpic 多方：Polytropic 准静态：Quasistatic

\documentclass{book}
%\documentclass[UTF8]{ctexbook}

%TODO——宏包
%页面格式
\usepackage{geometry}
	\geometry{a4paper, twoside, left = 2.54 cm, right = 2.54 cm, top = 3.18 cm, bottom = 3.18 cm, headheight = 3 cm}
\usepackage{fancyhdr}
	\fancyhf{}
	\fancyhead[EL]{\nouppercase{\CJKfamily{楷体} \leftmark}}
	\fancyhead[OR]{\nouppercase{\CJKfamily{楷体} \rightmark}}
	\fancyfoot[C]{\thepage}
	\renewcommand{\headrulewidth}{0 pt}

\usepackage{titlesec}
%	\newcommand{\sectionbreak}{\clearpage}

%交叉引用、超链接、PDF书签
\usepackage[hyperfootnotes = false]{hyperref}
	\hypersetup{
		bookmarksopen = true,
		bookmarksopenlevel = 1,
		pdfborder = 0 0 0
	}
%FIXME：防止PDF书签报错（LaTeX入门 p173）
%%使用 \texorpdfstring{TeX符号}{PDF书签符号}
%%例：非线性\texorpdfstring{$\sigma$}{σ}模型

%脚注
\usepackage[stable, perpage]{footmisc}
\usepackage{pifont}
	\renewcommand{\thefootnote}{\ding{\numexpr171+\value{footnote}}}
%\interfootnotelinepenalty = 10000

%数学支持
\usepackage{amsmath}
	\numberwithin{equation}{section}
	\renewcommand{\theequation}{\thesection.\arabic{equation}} %公式编号样式x.x.x
\usepackage{mathtools}
\usepackage[adobe-garamond]{mathdesign}
%\usepackage[urw-garamond]{mathdesign} %免费字体
\usepackage[ntheorem]{empheq}
\usepackage{bm}
%\usepackage{pifont} %已在“脚注”中调用
	\newcommand{\cmark}{\ding{51}}
	\newcommand{\xmark}{\ding{55}}

%字体
\usepackage[no-math]{fontspec}
	\setmainfont[
		Ligatures = TeX,
		BoldFont = AGaramondPro-Semibold.otf,
%		BoldFont = AGaramondPro-Bold.otf,
		ItalicFont = AGaramondPro-Italic.otf
	]{AGaramondPro-Regular.otf}
%	\setsansfont[
%		Ligatures = TeX,
%	]{Source Sans Pro}
	\setsansfont[
		Ligatures = TeX,
		ItalicFont = MyriadPro-It.otf
	]{MyriadPro-Regular.otf}
%	\newfontfamily{\helvetica}[
%		Ligatures = TeX,
%		BoldFont = HelveticaNeueLTPro-Roman.otf,
%		ItalicFont = HelveticaNeueLTPro-Roman.otf
%	]{HelveticaNeueLTPro-Roman.otf}
%	\newfontfamily{\myriad}[
%		Ligatures = TeX,
%		ItalicFont = MyriadPro-It.otf
%	]{MyriadPro-Regular.otf}

%中文支持
\usepackage[UTF8, heading = true]{ctex} %必须放在 xeCJK 之前调用，否则会覆盖字体设置
\usepackage{xeCJK}
	\xeCJKsetup{
		AutoFakeBold = false,
		AutoFakeSlant = false
	}
	\setCJKmainfont[
		BoldFont = 华文中宋,
		ItalicFont = 方正楷体_GBK,
		Mapping = fullwidth-stop
	]{方正书宋_GBK}
	\setCJKfamilyfont{中宋}[BoldFont = 华文中宋, ItalicFont = 华文中宋, Mapping = fullwidth-stop]{华文中宋}
	\setCJKfamilyfont{楷体}[BoldFont = 方正楷体_GBK, ItalicFont = 方正楷体_GBK, Mapping = fullwidth-stop]{方正楷体_GBK}
	\setCJKfamilyfont{仿宋}[BoldFont = 方正仿宋_GBK, ItalicFont = 方正仿宋_GBK, Mapping = fullwidth-stop]{方正仿宋_GBK}
	\setCJKfamilyfont{黑体}[BoldFont = 方正黑体_GBK, ItalicFont = 方正黑体_GBK, Mapping = fullwidth-stop]{方正黑体_GBK}
	\newcommand{\myHeavy}{\sffamily \CJKfamily{黑体}}
%\usepackage[UTF8, heading = true]{ctex}
	\pagestyle{fancy}
	\setcounter{secnumdepth}{4}
	\ctexset{
		part = {
			format = {\bfseries \Huge \centering},
			name = {第, 篇}
		},
		chapter = {
			format = {\bfseries \LARGE \raggedright},
		},
		section = {
			format = {\bfseries \Large \centering},
		},
		subsection = {
			format = {\bfseries},
			name = {,、\hspace{-1 em}},
			numbering = true,
			number = \chinese{subsection},
		}
	}

%符号支持
%\usepackage{gensymb}

%单位
\usepackage{siunitx}
	\sisetup{
		number-math-rm = \ensuremath,
		inter-unit-product = \ensuremath{{}\cdot{}},
		group-digits = integer,
		group-minimum-digits = 4,
		group-separator = \text{~}
	}

%项目列表
\usepackage{enumitem}
	\SetLabelAlign{leftalignwithindent}{\hspace{2.1 em} \makebox[1.6 em][l]{#1}}

%定理类环境
\usepackage[thmmarks, amsmath]{ntheorem}
	\theoremstyle{empty} %不带编号
	\theoremsymbol{\ensuremath{\Box}}
		\newtheorem{_myProof}{}
	
	\theoremstyle{empty} %不带编号
	\theoremheaderfont{\myHeavy}
	\theorembodyfont{\CJKfamily{楷体}}
	\theoremsymbol{}
%	\theoremindent \parindent
		\newtheorem{_myThm}{}
	
	\theoremstyle{plain} %带编号
	\theoremheaderfont{\myHeavy}
	\theorembodyfont{\normalfont}
	\theoremsymbol{}
%	\theoremsymbol{\ensuremath{\triangleleft}}
		\newtheorem{myExample}{例}[chapter]

%表格
\usepackage{array}
	\newcolumntype{M}{>{$}c<{$}} %数学模式，居中
\usepackage{tabularx}
	\newcolumntype{Y}{>{\centering\arraybackslash}X} %定宽居中
\usepackage{booktabs} %三线表

%图片
\usepackage{graphicx}
\usepackage{asymptote}
	\begin{asydef}
		usepackage("amsmath");
		usepackage("mathtools");
		usepackage("mathdesign", "adobe-garamond");
		usepackage("bm");
		usepackage("xeCJK");
			texpreamble("\setCJKmainfont[BoldFont = 华文中宋, ItalicFont = 方正楷体_GBK, Mapping = fullwidth-stop ]{方正书宋_GBK}");
			texpreamble("\setmainfont{AGaramondPro-Regular.otf}");
		usepackage("siunitx");
			texpreamble("\sisetup{
				number-math-rm = \ensuremath,
				inter-unit-product = \ensuremath{{}\cdot{}},
				group-digits = integer,
				group-minimum-digits = 4,
				group-separator = \text{~}
			}");
		
		size(6 cm);
		
		defaultpen(fontsize(9 pt));
		
		pen color1 = rgb(0.368417, 0.506779, 0.709798);
		pen color2 = rgb(0.880722, 0.611041, 0.142051);
		pen color3 = rgb(0.560181, 0.691569, 0.194885);
		pen color4 = rgb(0.922526, 0.385626, 0.209179);
		pen color5 = rgb(0.647624, 0.378160, 0.614037);
		pen color6 = rgb(0.772079, 0.431554, 0.102387);
		pen color7 = rgb(0.363898, 0.618501, 0.782349);
		pen color8 = rgb(0.972829, 0.621644, 0.073362);
	\end{asydef}

%浮动体、图标标题
\renewcommand{\thefigure}{\arabic{chapter}.\arabic{figure}}
\usepackage{floatrow}
	\floatsetup[table]{capposition = top}
\usepackage{caption}
	\DeclareCaptionFont{myCationFont}{\CJKfamily{楷体} \small}
	\captionsetup[figure]{
		font = myCationFont,
		labelsep = quad,
		skip = 10 pt,
		position = bottom
	}
	\captionsetup[table]{
		font = myCationFont,
		labelsep = quad,
		skip = 10 pt,
		position = top
	}


%TODO——自定义环境
%自定义定理（定律）
\newenvironment{myThm}[1]
	{\begin{_myThm}[\hskip 2 em #1]}
	{\end{_myThm}}
\newenvironment{myThm*}
	{\begin{_myThm} \hskip 2 em}
	{\end{_myThm}}
%自定义证明（说明）
\newenvironment{myProof}
	{\begin{_myProof} \normalfont \small \CJKfamily{仿宋} \hskip 2 em}
	{\end{_myProof}}
%自定义列表
\newenvironment{myEnum1}
	{\begin{enumerate} [label = \alph*), align = leftalignwithindent, leftmargin = 0 pt, itemsep = 10 pt, parsep = 0 pt, listparindent = 2 em]}
	{\end{enumerate}}
\newenvironment{myEnum2}
	{\begin{enumerate} [label = (\Roman*), align = leftalignwithindent, leftmargin = 2 em, topsep = 0pt, itemsep = 0 pt, parsep = 0 pt, listparindent = 2 em]}
	{\end{enumerate}}
%子公式
\newenvironment{mySubEq}
	{\subequations \renewcommand{\theequation}{\theparentequation-\alph{equation}}}
	{\endsubequations
	\ignorespacesafterend}
%带框公式
\newenvironment{boxedEq}
	{\empheq[box = \fbox]{equation}}
	{\endempheq}
%大括号公式
\newenvironment{braceEq}[1][align]
	{\mySubEq
		\setkeys{EmphEqEnv}{#1}
		\setkeys{EmphEqOpt}{left = \empheqlbrace}
		\EmphEqMainEnv}
	{\endEmphEqMainEnv \endmySubEq}
%自定义表格 {format}{caption}{label}
\newenvironment{myTable}[4][htb]
	{\begin{table}[#1] \centering \small \caption{#3} \label{#4} \begin{tabular}{#2}}
	{\end{tabular} \end{table}}

%TODO——自定义命令
%%大写正体字母
\newcommand{\rmY}{\mathrm{Y}}

%大写手写体字母
\newcommand{\scF}{\mathcal{F}}
\newcommand{\scH}{\mathcal{H}}
\newcommand{\scI}{\mathcal{I}}
\newcommand{\scJ}{\mathcal{J}}
\newcommand{\scL}{\mathcal{L}}
\newcommand{\scP}{\mathcal{P}}
\newcommand{\scR}{\mathcal{R}}
\newcommand{\scS}{\mathcal{S}}

%小写希腊字母
\renewcommand{\a}{\alpha}
\renewcommand{\b}{\beta}
\newcommand{\g}{\gamma}
\renewcommand{\d}{\delta}
\newcommand{\e}{\epsilon}
\newcommand{\ve}{\varepsilon}
\newcommand{\z}{\zeta}
\newcommand{\h}{\eta}
\newcommand{\q}{\theta}
\renewcommand{\th}{\theta}
\newcommand{\vq}{\vartheta}
\newcommand{\vth}{\vartheta}
\renewcommand{\i}{\iota}
\renewcommand{\k}{\varkappa} %这两个是相反的
\newcommand{\vk}{\kappa} %这两个是相反的
\renewcommand{\l}{\lambda}
\renewcommand{\m}{\mu}
\newcommand{\n}{\nu}
\newcommand{\x}{\xi}
%\renewcommand{\p}{\pi} %\p 已定义为正体π
\newcommand{\vp}{\varpi}
\renewcommand{\r}{\rho}
\newcommand{\vr}{\varrho}
\newcommand{\s}{\sigma}
\newcommand{\vs}{\varsigma}
\renewcommand{\t}{\tau}
\renewcommand{\u}{\upsilon}
\newcommand{\f}{\phi}
\renewcommand{\j}{\varphi}
\newcommand{\vf}{\varphi}
\renewcommand{\c}{\chi}
\newcommand{\y}{\psi}
\renewcommand{\o}{\omega}
\newcommand{\w}{\omega}

%大写希腊字母
\renewcommand{\G}{\varGamma}
\newcommand{\D}{\varDelta}
\newcommand{\Q}{\varTheta}
\newcommand{\Th}{\varTheta}
\renewcommand{\L}{\varLambda}
\newcommand{\X}{\varXi}
\renewcommand{\P}{\varPi}
\let \secSymbol = \S
\renewcommand{\S}{\varSigma}
\renewcommand{\U}{\varUpsilon}
\newcommand{\F}{\varPhi}
\newcommand{\Y}{\varPsi}
\renewcommand{\O}{\varOmega}
\newcommand{\W}{\varOmega}

%常数
\newcommand{\p}{\piup} %圆周率
\newcommand{\ee}{\mathrm{e}} %自然对数基
\newcommand{\ii}{\mathrm{i}} %虚数单位
\newcommand{\hb}{\hbar} %约化Planck常数
\newcommand{\kB}{k_{\mathrm{B}}} %Boltzmann常数
\newcommand{\pfs}{\e_0} %真空介电常数 Permittivity of Free Space

%数学修饰符号
\newcommand{\vecb}[1]{\bm{#1}} %矢量用斜体，根据ISO 80000-2/2009
\renewcommand{\bar}[1]{\overline{#1}} %上横线
\newcommand{\tl}[1]{\widetilde{#1}} %弯号
\newcommand{\op}[1]{\hat{#1}} %量子力学算符

%数学函数
\newcommand{\abs}[1]{\left\lvert #1 \right\rvert} %绝对值

%算子、算符
%已经调整好前后间距
\DeclareMathOperator{\incr}{\Delta\!}
\DeclareMathOperator{\dd}{\mathrm{d}\!} %微分 Differential D
\DeclareMathOperator{\db}{{\mkern2.5mu\mathchar'26\mkern-11mu\mathrm{d}}\!} %不完全微分 Inexact Differential D; Only for "Adobe Garamond" font!
\DeclareMathOperator{\pd}{\partial\!} %偏微分 Partial D
\DeclareMathOperator{\vd}{\deltaup\!} %变分 Variational D
\DeclareMathOperator{\laplac}{\nabla^2\!} %Laplace算子

%复合数学符号
\newcommand{\infint}{\int_{-\infty}^{+\infty}} %上下限无穷的积分
\newcommand{\iTonSum}{\sum_{i = 1}^{n}} %求和 i=1..n
\newcommand{\iToNSum}{\sum_{i = 1}^{N}} %求和 i=1..N
\newcommand{\nEnum}[2]{#1 = 1, \, 2, \, \dots, \, #2} %如x=1,2,...,n
%\newcommand{\coordinateLabel}[2]{#1\text{-}#2} %如p-V坐标
\newcommand{\myPartial}[3]{\left( \frac{\pd #1}{\pd #2} \right)_{#3}} %带括号的偏微分
\newcommand{\myPartialDisplay}[3]{\left( \dfrac{\pd #1}{\pd #2} \right)_{#3}}

%重定义符号
\newcommand{\approach}{\rightarrow} %趋向于
\newcommand{\eqdef}{\equiv} %定义等号（三线）
\newcommand{\dotTimes}{\cdot} %点乘

%公式用标点、文字
\newcommand{\comma}{\text{，}}
\newcommand{\fullstop}{\text{．}}
\newcommand{\semicomma}{\text{；}}
\newcommand{\const}{\text{const.}}

%正文用符号
%\catcode`\。 = \active
%\newcommand{。}{．} %句号用句点代替，已用xeCJK的Mapping代替
%\newcommand{\mySec}{{\fontspec{Times New Roman} \S \,}} %TODO：20160302 冲突；章节符号，Garamond 字体不好看
%%%%%%%%%%%%%%%%%%%%%%%%%%%%%%
%大写正体字母
\newcommand{\rmY}{\mathrm{Y}}

%大写手写体字母
\newcommand{\scF}{\mathcal{F}}
\newcommand{\scH}{\mathcal{H}}
\newcommand{\scI}{\mathcal{I}}
\newcommand{\scJ}{\mathcal{J}}
\newcommand{\scL}{\mathcal{L}}
\newcommand{\scP}{\mathcal{P}}
\newcommand{\scR}{\mathcal{R}}
\newcommand{\scS}{\mathcal{S}}

%小写希腊字母
\renewcommand{\a}{\alpha}
\renewcommand{\b}{\beta}
\newcommand{\g}{\gamma}
\renewcommand{\d}{\delta}
\newcommand{\e}{\epsilon}
\newcommand{\ve}{\varepsilon}
\newcommand{\z}{\zeta}
\newcommand{\h}{\eta}
\newcommand{\q}{\theta}
\renewcommand{\th}{\theta}
\newcommand{\vq}{\vartheta}
\newcommand{\vth}{\vartheta}
\renewcommand{\i}{\iota}
\renewcommand{\k}{\varkappa} %这两个是相反的
\newcommand{\vk}{\kappa} %这两个是相反的
\renewcommand{\l}{\lambda}
\renewcommand{\m}{\mu}
\newcommand{\n}{\nu}
\newcommand{\x}{\xi}
%\renewcommand{\p}{\pi} %\p 已定义为正体π
\newcommand{\vp}{\varpi}
\renewcommand{\r}{\rho}
\newcommand{\vr}{\varrho}
\newcommand{\s}{\sigma}
\newcommand{\vs}{\varsigma}
\renewcommand{\t}{\tau}
\renewcommand{\u}{\upsilon}
\newcommand{\f}{\phi}
\renewcommand{\j}{\varphi}
\newcommand{\vf}{\varphi}
\renewcommand{\c}{\chi}
\newcommand{\y}{\psi}
\renewcommand{\o}{\omega}
\newcommand{\w}{\omega}

%大写希腊字母
\renewcommand{\G}{\varGamma}
\newcommand{\D}{\varDelta}
\newcommand{\Q}{\varTheta}
\newcommand{\Th}{\varTheta}
\renewcommand{\L}{\varLambda}
\newcommand{\X}{\varXi}
\renewcommand{\P}{\varPi}
\let \secSymbol = \S
\renewcommand{\S}{\varSigma}
\renewcommand{\U}{\varUpsilon}
\newcommand{\F}{\varPhi}
\newcommand{\Y}{\varPsi}
\renewcommand{\O}{\varOmega}
\newcommand{\W}{\varOmega}

%常数
\newcommand{\p}{\piup} %圆周率
\newcommand{\ee}{\mathrm{e}} %自然对数基
\newcommand{\ii}{\mathrm{i}} %虚数单位
\newcommand{\hb}{\hbar} %约化Planck常数
\newcommand{\kB}{k_{\mathrm{B}}} %Boltzmann常数
\newcommand{\pfs}{\e_0} %真空介电常数 Permittivity of Free Space

%数学修饰符号
\newcommand{\vecb}[1]{\bm{#1}} %矢量用斜体，根据ISO 80000-2/2009
\renewcommand{\bar}[1]{\overline{#1}} %上横线
\newcommand{\tl}[1]{\widetilde{#1}} %弯号
\newcommand{\op}[1]{\hat{#1}} %量子力学算符

%数学函数
\newcommand{\abs}[1]{\left\lvert #1 \right\rvert} %绝对值

%算子、算符
%已经调整好前后间距
\DeclareMathOperator{\incr}{\Delta\!}
\DeclareMathOperator{\dd}{\mathrm{d}\!} %微分 Differential D
\DeclareMathOperator{\db}{{\mkern2.5mu\mathchar'26\mkern-11mu\mathrm{d}}\!} %不完全微分 Inexact Differential D; Only for "Adobe Garamond" font!
\DeclareMathOperator{\pd}{\partial\!} %偏微分 Partial D
\DeclareMathOperator{\vd}{\deltaup\!} %变分 Variational D
\DeclareMathOperator{\laplac}{\nabla^2\!} %Laplace算子

%复合数学符号
\newcommand{\infint}{\int_{-\infty}^{+\infty}} %上下限无穷的积分
\newcommand{\iTonSum}{\sum_{i = 1}^{n}} %求和 i=1..n
\newcommand{\iToNSum}{\sum_{i = 1}^{N}} %求和 i=1..N
\newcommand{\nEnum}[2]{#1 = 1, \, 2, \, \dots, \, #2} %如x=1,2,...,n
%\newcommand{\coordinateLabel}[2]{#1\text{-}#2} %如p-V坐标
\newcommand{\myPartial}[3]{\left( \frac{\pd #1}{\pd #2} \right)_{#3}} %带括号的偏微分
\newcommand{\myPartialDisplay}[3]{\left( \dfrac{\pd #1}{\pd #2} \right)_{#3}}

%重定义符号
\newcommand{\approach}{\rightarrow} %趋向于
\newcommand{\eqdef}{\equiv} %定义等号（三线）
\newcommand{\dotTimes}{\cdot} %点乘

%公式用标点、文字
\newcommand{\comma}{\text{，}}
\newcommand{\fullstop}{\text{．}}
\newcommand{\semicomma}{\text{；}}
\newcommand{\const}{\text{const.}}

%正文用符号
%\catcode`\。 = \active
%\newcommand{。}{．} %句号用句点代替，已用xeCJK的Mapping代替
%\newcommand{\mySec}{{\fontspec{Times New Roman} \S \,}} %TODO：20160302 冲突；章节符号，Garamond 字体不好看

%%%%%%%%%%%%%%%%%%%%%%%%%%%%%%

%强调
\newcommand{\emphA}[1]{{\myHeavy #1}}
\newcommand{\emphB}[1]{{\itshape #1}}

%引入文件
\let \oldInclude = \include
\renewcommand{\include}[1]{{\let \clearpage = \relax \oldInclude{#1}}}

%空行
\newcommand{\blankline}{\mbox{}}

%自定义列表项目（不可删除换行！）
\newcommand{\myItem}[1]{
	\item
	{\bfseries #1}
	
	
}

%引用格式
\newcommand{\secref}[1]{\secSymbol \ref{#1} }
\newcommand{\subsecref}[1]{第\ref{#1}小节}
\newcommand{\egref}[1]{例 \ref{#1}}

%公式标注
\newcommand{\myTag}[1]{\tag*{\CJKfamily{楷体} [#1]}}
\newcommand{\myTagNumbering}[1]{\refstepcounter{equation} \tag*{\CJKfamily{楷体} [#1] \, (\theequation)} }

%TODO——标题页
\title{
	\vspace{-4 cm}
	\bfseries
	热力学与统计物理I
}
\author{
	\CJKfamily{楷体}
	复旦大学\phantom{空格}陈焱
}
\date{
	\CJKfamily{楷体}
	\today
}

\begin{document}
%	\maketitle
%	
%	\frontmatter
%	{\let \cleardoublepage = \clearpage
%		\chapter{绪论}
%			本课程包含两部分内容：热力学、统计物理。

\emphA{热力学（thermodynamics）}是一种自上而下（top-down）的研究方式，它是形而上的、唯象的。用一句话可以形容热力学：“知其然而不知其所以然”。它主要研究宏观物理量之间的关系。

对热力学有重大贡献的物理学家有Carnot、Joule、Clausius、Kelvin等。

\blankline

\emphA{统计物理（statistical Physics）}是一种自下而上（bottom-up）的研究方式，它是形而下的、微观的。

对统计物理有重大贡献的物理学家见表 \ref{TAB_PHYSICIST_IN_STATISTICAL_PHYSICS}。

\begin{table}[htb]
	\begin{tabular}{c|c}
		\hline
		\emphA{阶段} & \emphA{人物} \\
		\hline
		经典统计 & Maxwell、Boltzmann、Gibbs、Einstein等 \\
		量子概念 & Planck、Einstein、Fermi、Dirac、Pauli、Bose等 \\
		量子统计 & von Neumann、Landau、Kramers、Pauli等 \\
		\hline
	\end{tabular}
	\caption{对统计物理有重大贡献的物理学家}
	\label{TAB_PHYSICIST_IN_STATISTICAL_PHYSICS}
\end{table}
统计物理可分为两个阶段。1860年至1902年，人们主要研究近独立子体系（简单地说，就是理想气体）；1902年以后，出现系综理论，开始研究凝聚态系统。

统计物理的基础可以概括为\emphB{等概率原理}。
%			
%		\chapter{符号表}
%			\include{Chapters/Front_Symbols}
%			
%		\chapter{缩略词表}
%			\begin{table}[h]
	\begin{tabularx}{\textwidth}{p{8 em}XX}
		ext			&	external			&	外界			\\
		max			&	maximum				&	最大值			\\
		min			&	minimum				&	最小值			\\
	\end{tabularx}
\end{table}

%	}
%	\mainmatter
%	\part{热力学}
		\chapter{热力学基础}
			% Chapter 1

\chapter{Chapter 1} % Main chapter title

\label{Chapter 1} % For referencing the chapter elsewhere, use \ref{Chapter1} 

%----------------------------------------------------------------------------------------

% Define some commands to keep the formatting separated from the content 
\newcommand{\keyword}[1]{\textbf{#1}}
\newcommand{\tabhead}[1]{\textbf{#1}}
\newcommand{\code}[1]{\texttt{#1}}
\newcommand{\file}[1]{\texttt{\bfseries#1}}
\newcommand{\option}[1]{\texttt{\itshape#1}}

%----------------------------------------------------------------------------------------

\section{Welcome and Thank You}
Welcome to this \LaTeX{} Thesis Template, a beautiful and easy to use template for writing a thesis using the \LaTeX{} typesetting system.

If you are writing a thesis (or will be in the future) and its subject is technical or mathematical (though it doesn't have to be), then creating it in \LaTeX{} is highly recommended as a way to make sure you can just get down to the essential writing without having to worry over formatting or wasting time arguing with your word processor.

\LaTeX{} is easily able to professionally typeset documents that run to hundreds or thousands of pages long. With simple mark-up commands, it automatically sets out the table of contents, margins, page headers and footers and keeps the formatting consistent and beautiful. One of its main strengths is the way it can easily typeset mathematics, even \option{heavy} mathematics. Even if those equations are the most horribly twisted and most difficult mathematical problems that can only be solved on a super-computer, you can at least count on \LaTeX{} to make them look stunning.

%----------------------------------------------------------------------------------------

\section{Learning \LaTeX{}}

\LaTeX{} is not a \textsc{wysiwyg} (What You See is What You Get) program, unlike word processors such as Microsoft Word or Apple's Pages. Instead, a document written for \LaTeX{} is actually a simple, plain text file that contains \emph{no formatting}. You tell \LaTeX{} how you want the formatting in the finished document by writing in simple commands amongst the text, for example, if I want to use \emph{italic text for emphasis}, I write the \verb|\emph{text}| command and put the text I want in italics in between the curly braces. This means that \LaTeX{} is a \enquote{mark-up} language, very much like HTML.

\subsection{A (not so short) Introduction to \LaTeX{}}

If you are new to \LaTeX{}, there is a very good eBook -- freely available online as a PDF file -- called, \enquote{The Not So Short Introduction to \LaTeX{}}. The book's title is typically shortened to just \emph{lshort}. You can download the latest version (as it is occasionally updated) from here:
\url{http://www.ctan.org/tex-archive/info/lshort/english/lshort.pdf}

It is also available in several other languages. Find yours from the list on this page: \url{http://www.ctan.org/tex-archive/info/lshort/}

It is recommended to take a little time out to learn how to use \LaTeX{} by creating several, small `test' documents, or having a close look at several templates on:\\ 
\url{http://www.LaTeXTemplates.com}\\ 
Making the effort now means you're not stuck learning the system when what you \emph{really} need to be doing is writing your thesis.

\subsection{A Short Math Guide for \LaTeX{}}

If you are writing a technical or mathematical thesis, then you may want to read the document by the AMS (American Mathematical Society) called, \enquote{A Short Math Guide for \LaTeX{}}. It can be found online here:
\url{http://www.ams.org/tex/amslatex.html}
under the \enquote{Additional Documentation} section towards the bottom of the page.

\subsection{Common \LaTeX{} Math Symbols}
There are a multitude of mathematical symbols available for \LaTeX{} and it would take a great effort to learn the commands for them all. The most common ones you are likely to use are shown on this page:
\url{http://www.sunilpatel.co.uk/latex-type/latex-math-symbols/}

You can use this page as a reference or crib sheet, the symbols are rendered as large, high quality images so you can quickly find the \LaTeX{} command for the symbol you need.

\subsection{\LaTeX{} on a Mac}
 
The \LaTeX{} distribution is available for many systems including Windows, Linux and Mac OS X. The package for OS X is called MacTeX and it contains all the applications you need -- bundled together and pre-customized -- for a fully working \LaTeX{} environment and work flow.
 
MacTeX includes a custom dedicated \LaTeX{} editor called TeXShop for writing your `\file{.tex}' files and BibDesk: a program to manage your references and create your bibliography section just as easily as managing songs and creating playlists in iTunes.

%----------------------------------------------------------------------------------------

\section{Getting Started with this Template}

If you are familiar with \LaTeX{}, then you should explore the directory structure of the template and then proceed to place your own information into the \emph{THESIS INFORMATION} block of the \file{main.tex} file. You can then modify the rest of this file to your unique specifications based on your degree/university. Section \ref{FillingFile} on page \pageref{FillingFile} will help you do this. Make sure you also read section \ref{ThesisConventions} about thesis conventions to get the most out of this template.

If you are new to \LaTeX{} it is recommended that you carry on reading through the rest of the information in this document.

Before you begin using this template you should ensure that its style complies with the thesis style guidelines imposed by your institution. In most cases this template style and layout will be suitable. If it is not, it may only require a small change to bring the template in line with your institution's recommendations. These modifications will need to be done on the \file{MastersDoctoralThesis.cls} file.

\subsection{About this Template}

This \LaTeX{} Thesis Template is originally based and created around a \LaTeX{} style file created by Steve R.\ Gunn from the University of Southampton (UK), department of Electronics and Computer Science. You can find his original thesis style file at his site, here:
\url{http://www.ecs.soton.ac.uk/~srg/softwaretools/document/templates/}

Steve's \file{ecsthesis.cls} was then taken by Sunil Patel who modified it by creating a skeleton framework and folder structure to place the thesis files in. The resulting template can be found on Sunil's site here:
\url{http://www.sunilpatel.co.uk/thesis-template}

Sunil's template was made available through \url{http://www.LaTeXTemplates.com} where it was modified many times based on user requests and questions. Version 2.0 and onwards of this template represents a major modification to Sunil's template and is, in fact, hardly recognisable. The work to make version 2.0 possible was carried out by \href{mailto:vel@latextemplates.com}{Vel} and Johannes Böttcher.

%----------------------------------------------------------------------------------------

\section{What this Template Includes}

\subsection{Folders}

This template comes as a single zip file that expands out to several files and folders. The folder names are mostly self-explanatory:

\keyword{Appendices} -- this is the folder where you put the appendices. Each appendix should go into its own separate \file{.tex} file. An example and template are included in the directory.

\keyword{Chapters} -- this is the folder where you put the thesis chapters. A thesis usually has about six chapters, though there is no hard rule on this. Each chapter should go in its own separate \file{.tex} file and they can be split as:
\begin{itemize}
\item Chapter 1: Introduction to the thesis topic
\item Chapter 2: Background information and theory
\item Chapter 3: (Laboratory) experimental setup
\item Chapter 4: Details of experiment 1
\item Chapter 5: Details of experiment 2
\item Chapter 6: Discussion of the experimental results
\item Chapter 7: Conclusion and future directions
\end{itemize}
This chapter layout is specialised for the experimental sciences, your discipline may be different.

\keyword{Figures} -- this folder contains all figures for the thesis. These are the final images that will go into the thesis document.

\subsection{Files}

Included are also several files, most of them are plain text and you can see their contents in a text editor. After initial compilation, you will see that more auxiliary files are created by \LaTeX{} or BibTeX and which you don't need to delete or worry about:

\keyword{example.bib} -- this is an important file that contains all the bibliographic information and references that you will be citing in the thesis for use with BibTeX. You can write it manually, but there are reference manager programs available that will create and manage it for you. Bibliographies in \LaTeX{} are a large subject and you may need to read about BibTeX before starting with this. Many modern reference managers will allow you to export your references in BibTeX format which greatly eases the amount of work you have to do.

\keyword{MastersDoctoralThesis.cls} -- this is an important file. It is the class file that tells \LaTeX{} how to format the thesis. 

\keyword{main.pdf} -- this is your beautifully typeset thesis (in the PDF file format) created by \LaTeX{}. It is supplied in the PDF with the template and after you compile the template you should get an identical version.

\keyword{main.tex} -- this is an important file. This is the file that you tell \LaTeX{} to compile to produce your thesis as a PDF file. It contains the framework and constructs that tell \LaTeX{} how to layout the thesis. It is heavily commented so you can read exactly what each line of code does and why it is there. After you put your own information into the \emph{THESIS INFORMATION} block -- you have now started your thesis!

Files that are \emph{not} included, but are created by \LaTeX{} as auxiliary files include:

\keyword{main.aux} -- this is an auxiliary file generated by \LaTeX{}, if it is deleted \LaTeX{} simply regenerates it when you run the main \file{.tex} file.

\keyword{main.bbl} -- this is an auxiliary file generated by BibTeX, if it is deleted, BibTeX simply regenerates it when you run the \file{main.aux} file. Whereas the \file{.bib} file contains all the references you have, this \file{.bbl} file contains the references you have actually cited in the thesis and is used to build the bibliography section of the thesis.

\keyword{main.blg} -- this is an auxiliary file generated by BibTeX, if it is deleted BibTeX simply regenerates it when you run the main \file{.aux} file.

\keyword{main.lof} -- this is an auxiliary file generated by \LaTeX{}, if it is deleted \LaTeX{} simply regenerates it when you run the main \file{.tex} file. It tells \LaTeX{} how to build the \emph{List of Figures} section.

\keyword{main.log} -- this is an auxiliary file generated by \LaTeX{}, if it is deleted \LaTeX{} simply regenerates it when you run the main \file{.tex} file. It contains messages from \LaTeX{}, if you receive errors and warnings from \LaTeX{}, they will be in this \file{.log} file.

\keyword{main.lot} -- this is an auxiliary file generated by \LaTeX{}, if it is deleted \LaTeX{} simply regenerates it when you run the main \file{.tex} file. It tells \LaTeX{} how to build the \emph{List of Tables} section.

\keyword{main.out} -- this is an auxiliary file generated by \LaTeX{}, if it is deleted \LaTeX{} simply regenerates it when you run the main \file{.tex} file.

So from this long list, only the files with the \file{.bib}, \file{.cls} and \file{.tex} extensions are the most important ones. The other auxiliary files can be ignored or deleted as \LaTeX{} and BibTeX will regenerate them.

%----------------------------------------------------------------------------------------

\section{Filling in Your Information in the \file{main.tex} File}\label{FillingFile}

You will need to personalise the thesis template and make it your own by filling in your own information. This is done by editing the \file{main.tex} file in a text editor or your favourite LaTeX environment.

Open the file and scroll down to the third large block titled \emph{THESIS INFORMATION} where you can see the entries for \emph{University Name}, \emph{Department Name}, etc \ldots

Fill out the information about yourself, your group and institution. You can also insert web links, if you do, make sure you use the full URL, including the \code{http://} for this. If you don't want these to be linked, simply remove the \verb|\href{url}{name}| and only leave the name.

When you have done this, save the file and recompile \code{main.tex}. All the information you filled in should now be in the PDF, complete with web links. You can now begin your thesis proper!

%----------------------------------------------------------------------------------------

\section{The \code{main.tex} File Explained}

The \file{main.tex} file contains the structure of the thesis. There are plenty of written comments that explain what pages, sections and formatting the \LaTeX{} code is creating. Each major document element is divided into commented blocks with titles in all capitals to make it obvious what the following bit of code is doing. Initially there seems to be a lot of \LaTeX{} code, but this is all formatting, and it has all been taken care of so you don't have to do it.

Begin by checking that your information on the title page is correct. For the thesis declaration, your institution may insist on something different than the text given. If this is the case, just replace what you see with what is required in the \emph{DECLARATION PAGE} block.

Then comes a page which contains a funny quote. You can put your own, or quote your favourite scientist, author, person, and so on. Make sure to put the name of the person who you took the quote from.

Following this is the abstract page which summarises your work in a condensed way and can almost be used as a standalone document to describe what you have done. The text you write will cause the heading to move up so don't worry about running out of space.

Next come the acknowledgements. On this page, write about all the people who you wish to thank (not forgetting parents, partners and your advisor/supervisor).

The contents pages, list of figures and tables are all taken care of for you and do not need to be manually created or edited. The next set of pages are more likely to be optional and can be deleted since they are for a more technical thesis: insert a list of abbreviations you have used in the thesis, then a list of the physical constants and numbers you refer to and finally, a list of mathematical symbols used in any formulae. Making the effort to fill these tables means the reader has a one-stop place to refer to instead of searching the internet and references to try and find out what you meant by certain abbreviations or symbols.

The list of symbols is split into the Roman and Greek alphabets. Whereas the abbreviations and symbols ought to be listed in alphabetical order (and this is \emph{not} done automatically for you) the list of physical constants should be grouped into similar themes.

The next page contains a one line dedication. Who will you dedicate your thesis to?

Finally, there is the block where the chapters are included. Uncomment the lines (delete the \code{\%} character) as you write the chapters. Each chapter should be written in its own file and put into the \emph{Chapters} folder and named \file{Chapter1}, \file{Chapter2}, etc\ldots Similarly for the appendices, uncomment the lines as you need them. Each appendix should go into its own file and placed in the \emph{Appendices} folder.

After the preamble, chapters and appendices finally comes the bibliography. The bibliography style (called \option{authoryear}) is used for the bibliography and is a fully featured style that will even include links to where the referenced paper can be found online. Do not underestimate how grateful your reader will be to find that a reference to a paper is just a click away. Of course, this relies on you putting the URL information into the BibTeX file in the first place.

%----------------------------------------------------------------------------------------

\section{Thesis Features and Conventions}\label{ThesisConventions}

To get the best out of this template, there are a few conventions that you may want to follow.

One of the most important (and most difficult) things to keep track of in such a long document as a thesis is consistency. Using certain conventions and ways of doing things (such as using a Todo list) makes the job easier. Of course, all of these are optional and you can adopt your own method.

\subsection{Printing Format}

This thesis template is designed for double sided printing (i.e. content on the front and back of pages) as most theses are printed and bound this way. Switching to one sided printing is as simple as uncommenting the \option{oneside} option of the \code{documentclass} command at the top of the \file{main.tex} file. You may then wish to adjust the margins to suit specifications from your institution.

The headers for the pages contain the page number on the outer side (so it is easy to flick through to the page you want) and the chapter name on the inner side.

The text is set to 11 point by default with single line spacing, again, you can tune the text size and spacing should you want or need to using the options at the very start of \file{main.tex}. The spacing can be changed similarly by replacing the \option{singlespacing} with \option{onehalfspacing} or \option{doublespacing}.

\subsection{Using US Letter Paper}

The paper size used in the template is A4, which is the standard size in Europe. If you are using this thesis template elsewhere and particularly in the United States, then you may have to change the A4 paper size to the US Letter size. This can be done in the margins settings section in \file{main.tex}.

Due to the differences in the paper size, the resulting margins may be different to what you like or require (as it is common for institutions to dictate certain margin sizes). If this is the case, then the margin sizes can be tweaked by modifying the values in the same block as where you set the paper size. Now your document should be set up for US Letter paper size with suitable margins.

\subsection{References}

The \code{biblatex} package is used to format the bibliography and inserts references such as this one \parencite{Reference1}. The options used in the \file{main.tex} file mean that the in-text citations of references are formatted with the author(s) listed with the date of the publication. Multiple references are separated by semicolons (e.g. \parencite{Reference2, Reference1}) and references with more than three authors only show the first author with \emph{et al.} indicating there are more authors (e.g. \parencite{Reference3}). This is done automatically for you. To see how you use references, have a look at the \file{Chapter1.tex} source file. Many reference managers allow you to simply drag the reference into the document as you type.

Scientific references should come \emph{before} the punctuation mark if there is one (such as a comma or period). The same goes for footnotes\footnote{Such as this footnote, here down at the bottom of the page.}. You can change this but the most important thing is to keep the convention consistent throughout the thesis. Footnotes themselves should be full, descriptive sentences (beginning with a capital letter and ending with a full stop). The APA6 states: \enquote{Footnote numbers should be superscripted, [...], following any punctuation mark except a dash.} The Chicago manual of style states: \enquote{A note number should be placed at the end of a sentence or clause. The number follows any punctuation mark except the dash, which it precedes. It follows a closing parenthesis.}

The bibliography is typeset with references listed in alphabetical order by the first author's last name. This is similar to the APA referencing style. To see how \LaTeX{} typesets the bibliography, have a look at the very end of this document (or just click on the reference number links in in-text citations).

\subsubsection{A Note on bibtex}

The bibtex backend used in the template by default does not correctly handle unicode character encoding (i.e. "international" characters). You may see a warning about this in the compilation log and, if your references contain unicode characters, they may not show up correctly or at all. The solution to this is to use the biber backend instead of the outdated bibtex backend. This is done by finding this in \file{main.tex}: \option{backend=bibtex} and changing it to \option{backend=biber}. You will then need to delete all auxiliary BibTeX files and navigate to the template directory in your terminal (command prompt). Once there, simply type \code{biber main} and biber will compile your bibliography. You can then compile \file{main.tex} as normal and your bibliography will be updated. An alternative is to set up your LaTeX editor to compile with biber instead of bibtex, see \href{http://tex.stackexchange.com/questions/154751/biblatex-with-biber-configuring-my-editor-to-avoid-undefined-citations/}{here} for how to do this for various editors.

\subsection{Tables}

Tables are an important way of displaying your results, below is an example table which was generated with this code:

{\small
\begin{verbatim}
\begin{table}
\caption{The effects of treatments X and Y on the four groups studied.}
\label{tab:treatments}
\centering
\begin{tabular}{l l l}
\toprule
\tabhead{Groups} & \tabhead{Treatment X} & \tabhead{Treatment Y} \\
\midrule
1 & 0.2 & 0.8\\
2 & 0.17 & 0.7\\
3 & 0.24 & 0.75\\
4 & 0.68 & 0.3\\
\bottomrule\\
\end{tabular}
\end{table}
\end{verbatim}
}

\begin{table}
\caption{The effects of treatments X and Y on the four groups studied.}
\label{tab:treatments}
\centering
\begin{tabular}{l l l}
\toprule
\tabhead{Groups} & \tabhead{Treatment X} & \tabhead{Treatment Y} \\
\midrule
1 & 0.2 & 0.8\\
2 & 0.17 & 0.7\\
3 & 0.24 & 0.75\\
4 & 0.68 & 0.3\\
\bottomrule\\
\end{tabular}
\end{table}

You can reference tables with \verb|\ref{<label>}| where the label is defined within the table environment. See \file{Chapter1.tex} for an example of the label and citation (e.g. Table~\ref{tab:treatments}).

\subsection{Figures}

There will hopefully be many figures in your thesis (that should be placed in the \emph{Figures} folder). The way to insert figures into your thesis is to use a code template like this:
\begin{verbatim}
\begin{figure}
\centering
\includegraphics{Figures/Electron}
\decoRule
\caption[An Electron]{An electron (artist's impression).}
\label{fig:Electron}
\end{figure}
\end{verbatim}
Also look in the source file. Putting this code into the source file produces the picture of the electron that you can see in the figure below.

\begin{figure}[th]
\centering
\includegraphics{Figures/Electron}
\decoRule
\caption[An Electron]{An electron (artist's impression).}
\label{fig:Electron}
\end{figure}

Sometimes figures don't always appear where you write them in the source. The placement depends on how much space there is on the page for the figure. Sometimes there is not enough room to fit a figure directly where it should go (in relation to the text) and so \LaTeX{} puts it at the top of the next page. Positioning figures is the job of \LaTeX{} and so you should only worry about making them look good!

Figures usually should have captions just in case you need to refer to them (such as in Figure~\ref{fig:Electron}). The \verb|\caption| command contains two parts, the first part, inside the square brackets is the title that will appear in the \emph{List of Figures}, and so should be short. The second part in the curly brackets should contain the longer and more descriptive caption text.

The \verb|\decoRule| command is optional and simply puts an aesthetic horizontal line below the image. If you do this for one image, do it for all of them.

\LaTeX{} is capable of using images in pdf, jpg and png format.

\subsection{Typesetting mathematics}

If your thesis is going to contain heavy mathematical content, be sure that \LaTeX{} will make it look beautiful, even though it won't be able to solve the equations for you.

The \enquote{Not So Short Introduction to \LaTeX} (available on \href{http://www.ctan.org/tex-archive/info/lshort/english/lshort.pdf}{CTAN}) should tell you everything you need to know for most cases of typesetting mathematics. If you need more information, a much more thorough mathematical guide is available from the AMS called, \enquote{A Short Math Guide to \LaTeX} and can be downloaded from:
\url{ftp://ftp.ams.org/pub/tex/doc/amsmath/short-math-guide.pdf}

There are many different \LaTeX{} symbols to remember, luckily you can find the most common symbols in \href{http://ctan.org/pkg/comprehensive}{The Comprehensive \LaTeX~Symbol List}.

You can write an equation, which is automatically given an equation number by \LaTeX{} like this:
\begin{verbatim}
\begin{equation}
E = mc^{2}
\label{eqn:Einstein}
\end{equation}
\end{verbatim}

This will produce Einstein's famous energy-matter equivalence equation:
\begin{equation}
E = mc^{2}
\label{eqn:Einstein}
\end{equation}

All equations you write (which are not in the middle of paragraph text) are automatically given equation numbers by \LaTeX{}. If you don't want a particular equation numbered, use the unnumbered form:
\begin{verbatim}
\[ a^{2}=4 \]
\end{verbatim}

%----------------------------------------------------------------------------------------

\section{Sectioning and Subsectioning}

You should break your thesis up into nice, bite-sized sections and subsections. \LaTeX{} automatically builds a table of Contents by looking at all the \verb|\chapter{}|, \verb|\section{}|  and \verb|\subsection{}| commands you write in the source.

The Table of Contents should only list the sections to three (3) levels. A \verb|chapter{}| is level zero (0). A \verb|\section{}| is level one (1) and so a \verb|\subsection{}| is level two (2). In your thesis it is likely that you will even use a \verb|subsubsection{}|, which is level three (3). The depth to which the Table of Contents is formatted is set within \file{MastersDoctoralThesis.cls}. If you need this changed, you can do it in \file{main.tex}.

%----------------------------------------------------------------------------------------

\section{In Closing}

You have reached the end of this mini-guide. You can now rename or overwrite this pdf file and begin writing your own \file{Chapter1.tex} and the rest of your thesis. The easy work of setting up the structure and framework has been taken care of for you. It's now your job to fill it out!

Good luck and have lots of fun!

\begin{flushright}
Guide written by ---\\
Sunil Patel: \href{http://www.sunilpatel.co.uk}{www.sunilpatel.co.uk}\\
Vel: \href{http://www.LaTeXTemplates.com}{LaTeXTemplates.com}
\end{flushright}

			
		\chapter{均匀系统的平衡特性}
			\section{Maxwell关系} \label{SEC_Maxwell关系}
	\subsection{热力学函数的全微分}
		\begin{myEnum1}
			\myItem{$U$ 的全微分}
				根据热力学基本微分方程［见式~\eqref{EQ_FUNDAMENTAL_EQUATION_FOR_PVT_SYSTEM}］
				\begin{equation} \label{EQ_dU=TdS-pdV_IN_MAXWELL_RELATION}
					\dd U = T \dd S + p \dd V \comma
				\end{equation}
				可知内能 $U$ 是 $S$ 和 $V$ 的函数，即
				\begin{equation}
					U = U(S, \, V) \fullstop
				\end{equation}
				因此
				\begin{equation}
					\dd U = \myPartial{U}{S}{V} \dd S + \myPartial{U}{V}{S} \dd V \fullstop
				\end{equation}
				与式~\eqref{EQ_dU=TdS-pdV_IN_MAXWELL_RELATION} 进行比较，可得
				\begin{braceEq}
					& T = \myPartial{U}{S}{V} \comma  \label{EQ_T_IN_PARTIAL_DIFFERENTIAL_OF_U} \\
					& p = -\myPartial{U}{V}{S}  \label{EQ_p_IN_PARTIAL_DIFFERENTIAL_OF_U} \fullstop
				\end{braceEq}
				$U$ 的二阶偏导数与其先后次序无关，即
				\begin{equation}
					\frac{\pd^{\:2} \! U}{\pd V \pd S} = \frac{\pd^{\:2} \! U}{\pd S \pd V}%HACK：20160330 偏微分空格
				\end{equation}
				或
				\begin{equation}
					\left[ \frac{\pd}{\pd V} \myPartial{U}{S}{V}\right]_S
					= \left[ \frac{\pd}{\pd S} \myPartial{U}{V}{S}\right]_V \fullstop
				\end{equation}
				利用该式，对式~\eqref{EQ_T_IN_PARTIAL_DIFFERENTIAL_OF_U} 和式~\eqref{EQ_p_IN_PARTIAL_DIFFERENTIAL_OF_U} 两边分别求导，可得
				\begin{equation} \label{EQ_MAXWELL_RELATION_FROM_U}
					\myPartial{T}{V}{S} = -\myPartial{p}{S}{V} \fullstop
				\end{equation}
				
			\myItem{$H$ 的全微分}
				因为
				\begin{equation}
					H \eqdef U + p V \comma
				\end{equation}
				所以
				\begin{align}
					\dd H &= T \dd S - p \dd V + \dd \; (p V) \notag \\
					&= T \dd S - p \dd V + (V \dd p + p \dd V) \notag \\
					&= T \dd S + V \dd p \fullstop
				\end{align}%HACK：20160330 微分d后接空格，间距有问题
				这说明焓 $H$ 是 $S$ 和 $p$ 的函数，即
				\begin{equation}
					H = H(S, \, p) \fullstop
				\end{equation}
				因此
				\begin{equation}
					\dd H = \myPartial{H}{S}{p} \dd S + \myPartial{H}{p}{S} \dd p \fullstop
				\end{equation}
				进而
				\begin{braceEq}
					& T = \myPartial{H}{S}{p} \comma \\
					& V = \myPartial{H}{p}{S} \fullstop
				\end{braceEq}
				与上文类似，对它们两边分别求导，并利用混合偏导数定理，就得到
				\begin{equation} \label{EQ_MAXWELL_RELATION_FROM_H}
					\myPartial{T}{p}{S} = \myPartial{V}{S}{p} \fullstop
				\end{equation}
				
				\myItem{$F$ 的全微分}
				因为
				\begin{equation}
					F \eqdef U - T S \comma
				\end{equation}
				所以
				\begin{equation} \label{EQ_dF=-SdT-pdV_IN_MAXWELL_RELATION}
					\dd F = -S \dd T - p \dd V \fullstop
				\end{equation}
				这说明自由能 $F$ 是 $T$ 和 $V$ 的函数，即
				\begin{equation}
					F = F(T, \, V) \fullstop
				\end{equation}
				因此
				\begin{equation}
					\dd F = \myPartial{F}{T}{V} \dd T + \myPartial{F}{V}{T} \dd V \fullstop
				\end{equation}
				进而
				\begin{braceEq}
					& S = -\myPartial{F}{T}{V} \comma \\
					& p = -\myPartial{F}{V}{T} \fullstop
				\end{braceEq}
				两边分别求导，得
				\begin{equation} \label{EQ_MAXWELL_RELATION_FROM_F}
					\myPartial{S}{V}{T} = \myPartial{p}{T}{V} \fullstop
				\end{equation}
			
			\myItem{$G$ 的全微分}
				因为
				\begin{equation}
					G \eqdef U - T S + p V \comma
				\end{equation}
				所以
				\begin{equation}
				\dd G = -S \dd T + V \dd p \fullstop
				\end{equation}
				这说明Gibbs自由能是 $T$ 和 $p$ 的函数，即
				\begin{equation}
					G = G(T, \, p) \fullstop
				\end{equation}
				因此
				\begin{equation}
					\dd H = \myPartial{G}{G}{p} \dd T + \myPartial{G}{p}{T} \dd T \fullstop
				\end{equation}
				进而
				\begin{braceEq}
					& S = -\myPartial{G}{T}{p} \comma \\
					& V = \myPartial{G}{p}{T} \fullstop
				\end{braceEq}
				两边分别求导，得
				\begin{equation} \label{EQ_MAXWELL_RELATION_FROM_G}
				\myPartial{S}{p}{T} = -\myPartial{V}{T}{p} \fullstop
				\end{equation}
		\end{myEnum1}
		
		\blankline
		
		式~\eqref{EQ_MAXWELL_RELATION_FROM_U}、\eqref{EQ_MAXWELL_RELATION_FROM_H}、\eqref{EQ_MAXWELL_RELATION_FROM_F} 和 \eqref{EQ_MAXWELL_RELATION_FROM_G} 称为\emphA{Maxwell关系}。
		
	\subsection{Legendre变换}
		\emphA{Legendre变换}是指等式
		\begin{equation}
			x \dd y = \dd \; (x y) - y \dd x \comma
		\end{equation}
		它把变量从 $x$ 变为了 $y$。
		
		在经典力学中，Legendre变换被用来从Lagrange表述导出Hamilton表述：
		\begin{equation}
			\scL (q_\a, \, \dot{q_\a}) \quad \rightarrow \quad \scH (q_\a, \, p_\a) = \sum_\a \dot{q_\a} p_\a - \scL (q_\a, \, \dot{q_\a}) \fullstop
		\end{equation}
		
		对式~\eqref{EQ_dU=TdS-pdV_IN_MAXWELL_RELATION} 应用Legendre变换
		\begin{equation}
			T \dd S = \dd \; (T S) - S \dd T
		\end{equation}
		并移项，就得到
		\begin{equation}
			\dd U - \dd \; (T S) = \dd \; (U - T S)= \dd F = -S \dd T + p \dd V \fullstop
		\end{equation}
		这也就是用\emphB{自由能}表示的热力学基本方程，即式~\eqref{EQ_dF=-SdT-pdV_IN_MAXWELL_RELATION}。由此，利用偏微分关系稍做计算，就得到了Maxwell关系 \eqref{EQ_MAXWELL_RELATION_FROM_F} 式。用同样的手法，也可以得到其他几个Maxwell关系。
		
		\begin{myTable}[H]{MM>{$}l<{$}}{Maxwell关系}{TAB_MAXWELL_RELATION}
			\toprule
			\text{\emphA{基本微分方程}} & \text{\emphA{自然变量}} & \text{\emphA{Maxwell关系}} \\
			\midrule
			\dd U = T \dd S - p \dd V & (S, \, V) & \myPartialDisplay{T}{V}{S} = -\myPartialDisplay{p}{S}{V} \\
			\dd H = \dd \; (U + p V) = T \dd S + V \dd p & (S, \, p) & \myPartialDisplay{T}{p}{S} = \myPartialDisplay{V}{S}{p} \\
			\dd F = \dd \; (U - T S) = -S \dd T - p \dd V & (T, \, V) & \myPartialDisplay{S}{V}{T} = \myPartialDisplay{p}{T}{V} \\
			\dd G = \dd \; (U - T S + p V) = -S \dd T + V \dd p & (T, \, p) & \myPartialDisplay{S}{p}{T} = -\myPartialDisplay{V}{T}{p} \\
			\bottomrule
		\end{myTable}%FIXME：20160330 表格格式
		
		表~\ref{TAB_MAXWELL_RELATION} 总结了上文的推导结果。其中的“\emphA{自然变量}”指经Legendre变换后的自由变量，它们直接出现在热力学基本方程中。
		
		热力学变量（或函数）［见 \secref{SEC_平衡态及其描述}\subsecref{SUBSEC_平衡态的描述}］可分为两种：\emphA{可测量量}和\emphA{不可测量量}。物态方程中的 $p$、$V$、$T$，与物态方程有关的 $\a$、$\b$、$\k_T$，以及热容，都是可测量量；而 $U$、$S$、$H$、$F$、$G$ 等状态函数，以及 $(\pd S / \pd V)_T$、$(\pd H / \pd p)_T$ 等偏微分，都是不可测量量。Maxwell关系的作用，就是把不可测量量用可测量量来表示。
		
	\subsection{简单应用} \label{SUBSEC_简单应用_OF_MAXWELL关系}
		\begin{myExample} \label{EG_C_p-C_V}
			计算 $C_p - C_V$。\footnote{
				注意与式~\eqref{EQ_C_p-C_V_FOR_IDEAL_GAS} 对比，那里要求理想气体。
			}%TODO：20160330 式1.4.8有问题
			
			根据 \secref{SEC_热容与焓；理想气体的性质}\subsecref{SUBSEC_热容与焓} 中的推导，
			\begin{braceEq}
				C_p &= \myPartial{H}{T}{p} \comma \\
				C_V &= \myPartial{U}{T}{V} \fullstop
			\end{braceEq}
			利用 $\db Q = T \dd S$，可得%TODO：20160330 为什么dQ=TdS(可逆？)
			\begin{braceEq}
				C_p &= T \myPartial{S}{T}{p} \comma \label{EQ_C_p_IN_T_AND_S} \\
				C_V &= T \myPartial{S}{T}{V} \fullstop \label{EQ_C_V_IN_T_AND_S} 
			\end{braceEq}%TODO：20160330 推导过程
			把 $S = S(T, \, p)$ 看成复合函数的形式，即 $S[T, \, V(T, \, p)]$，因此
			\begin{equation}
				\myPartial{S}{T}{p} = \myPartial{S}{T}{V} + \myPartial{S}{V}{T} \myPartial{V}{T}{p} \fullstop
			\end{equation}
			于是
			\begin{align}
				C_p - C_V &= T \left[ \myPartial{S}{T}{p} - \myPartial{S}{T}{V} \right] \notag \\
				&= T \myPartial{S}{V}{T} \myPartial{V}{T}{p} \notag \\%FIXME：20160401 强调
				&= T \myPartial{p}{T}{V} \myPartial{V}{T}{p} \fullstop
				\myTagNumbering{Maxwell关系}  \label{EQ_C_p-C_V_WITH_MAXWELL_RELATION_PART_1}
			\end{align}
			根据\emphB{偏导数三乘积法则} \eqref{EQ_CYCLIC_RELATION} 式，有
			\begin{equation}
				\myPartial{p}{T}{V} = -\myPartial{V}{T}{p} \myPartial{p}{V}{T} \semicomma
			\end{equation}
			根据\emphB{倒数关系} \eqref{EQ_RECIPROCITY_RELATION} 式，有
			\begin{equation}
				\myPartial{V}{T}{p} = \myPartial{T}{V}{p}^{-1} \fullstop
			\end{equation}
			%FIXME：20160711 交叉引用内容出错
			把以上两式代入式~\eqref{EQ_C_p-C_V_WITH_MAXWELL_RELATION_PART_1}，得
			\begin{align}
				C_p - C_V &= T \myPartial{p}{T}{V} \myPartial{V}{T}{p} \notag \\
				&= T \left[ -\myPartial{V}{T}{p} \myPartial{p}{V}{T} \right] \myPartial{V}{T}{p} \notag \\
				&= -T \myPartial{V}{T}{p}^2 \myPartial{V}{p}{T}^{-1} \fullstop \label{EQ_C_p-C_V_WITH_MAXWELL_RELATION_PART_2}
			\end{align}
			再把\emphB{膨胀系数}、\emphB{等温压缩系数}的定义
			\begin{braceEq}
				&\a = \frac{1}{V} \myPartial{V}{T}{p} \comma \\
				&\k_T = -\frac{1}{V} \myPartial{V}{p}{T}
			\end{braceEq}
			代入式~\eqref{EQ_C_p-C_V_WITH_MAXWELL_RELATION_PART_2}，便得到
			\begin{align}
				C_p - C_V &= -T (V \a)^2 \left( -\frac{1}{V \k_T} \right) \notag \\
				&= T V \frac{\a^2}{\k_T} \label{EQ_C_p-C_V_WITH_MAXWELL_RELATION_PART_3} \fullstop
			\end{align}
			这几个量都是可测量量。热力学稳定相中，有 $\k_T > 0$。因此 $C_p$ 始终大于 $C_V$。%FIXME：20160401 后续交叉引用
			
			对于理想气体，有
			\begin{braceEq}
				&\a = \frac{1}{V} \myPartial{V}{T}{p}
				= \frac{1}{V} \left( \frac{\pd}{\pd T} \frac{n R T}{p} \right)_p
				= \frac{n R}{p V} \comma \\
				&\k_T = -\frac{1}{V} \myPartial{V}{p}{T}
				= -\frac{1}{V} \left( \frac{\pd}{\pd p} \frac{n R T}{p} \right)_T
				= \frac{n R T}{p^2 V} \fullstop
			\end{braceEq}
			根据式~\eqref{EQ_C_p-C_V_WITH_MAXWELL_RELATION_PART_3}，可知
			\begin{equation}
				C_p - C_V = T V \left( \frac{n R}{p V} \right)^2 \left( \frac{p^2 V}{n R T} \right) = n R \comma
			\end{equation}
			这与式~\eqref{EQ_C_p-C_V_FOR_IDEAL_GAS} 是一致的。
		\end{myExample}
		
		\begin{myExample} \label{EG_pd_U/pd_V_WITH_FIXED_T}
			计算 $(\pd U/ \pd V)_T$。
			
			选取 $T$ 和 $V$ 作为独立变量，则
			\begin{align}
				\dd U &= T \dd S - p \dd V \notag \\
				&= T \left[ \myPartial{S}{T}{V} \dd T + \myPartial{S}{V}{T} \dd T \right] - p \dd V \notag \\
				&= T \myPartial{S}{T}{V} \dd T + \left[ T \myPartial{S}{V}{T} - p \right] \dd V \comma
			\end{align}
			因此
			\begin{align}
				\myPartial{U}{V}{T} &= T \myPartial{S}{V}{T} - p \notag \\
				&= T \myPartial{p}{T}{V} - p \fullstop \label{EQ_pd_U/pd_V_WITH_FIXED_T}
			\end{align}
			这里出现的都是状态变量或是与状态方程直接相关的量，因而也是可测量量。
			
			对于理想气体，$p V = n R T$，因此
			\begin{equation}
				\myPartial{p}{T}{V} = \left( \frac{\pd}{\pd T} \frac{n R T}{V} \right)_V = \frac{n R}{V} = \frac{p}{T} \comma
			\end{equation}
			代入式~\eqref{EQ_pd_U/pd_V_WITH_FIXED_T}，可以发现
			\begin{equation}
				\myPartial{U}{V}{T} = 0 \fullstop
			\end{equation}
			这说明 $U = U(T)$，即内能 $U$ 仅是温度 $T$ 的函数，而与体积 $V$ 无关。这就是\secref{SEC_热容与焓；理想气体的性质}\subsecref{SUBSEC_理想气体的性质}中提到的利用状态方程来证明内能只与温度有关的方法。
			
			对于van der Waals气体，其状态方程 ［式~\eqref{EQ_VAN_DER_WAALS_GAS_STATE_EQUATION}］为
			\begin{equation}
				\left( p + \frac{n^2 a}{V^2} \right) (V - n b) = n R T \fullstop
			\end{equation}
			在固定 $V$ 的条件下关于 $T$ 求偏导，得
			\begin{equation}
				\myPartial{p}{T}{V} (V - n b) + \left( p + \frac{n^2 a}{V^2} \right) \cdot 0 = n R \comma
			\end{equation}
			即
			\begin{equation}
				\myPartial{p}{T}{V} = \frac{n R}{V - n b} \fullstop
			\end{equation}
			因此
			\begin{equation}
				\myPartial{U}{V}{T} = \frac{n R T}{V - n b} - p = \frac{n^2 a}{V^2} \fullstop
			\end{equation}
			可见van der Waals气体的内能与体积\emphB{有关}。%TODO：20160401 与温度T有关？
		\end{myExample}
		
		\begin{myExample}
			证明绝热压缩系数与等温压缩系数之比
			\begin{equation}
				\frac{\k_\text{a}}{\k_T} = \frac{C_V}{C_p} \fullstop
			\end{equation}%HACK：20160401 绝热过程下标用\text{a}
			
			可逆的绝热过程是等熵过程，因此把 $\k_\text{a}$ 的下标改用“$S$”。%TODO：20160401 为什么可逆
			
			根据定义
			\begin{braceEq}
				\k_S &\eqdef -\frac{1}{V} \myPartial{V}{p}{S} \comma \\
				\k_T &\eqdef -\frac{1}{V} \myPartial{V}{p}{T} \fullstop
			\end{braceEq}
			因此
			\begin{align}
				\frac{\k_S}{\k_T} &= \left. \myPartialDisplay{V}{p}{S} \, \middle/ \, \myPartialDisplay{V}{p}{T} \right. \notag \\
				&= \left. \frac{ \pd \; (V, \, S) }{ \pd \; (p, \, S) } \, \middle/ \, \frac{ \pd \; (V, \, T) }{ \pd \; (p, \, T) } \right. \notag \\
				&= \left. \frac{ \pd \; (V, \, S) }{ \pd \; (V, \, T) } \, \middle/ \, \frac{ \pd \; (p, \, S) }{ \pd \; (p, \, T) } \right. \myTag{见式} \\%FIXME：20160402 交叉引用没写
				&= \left. \myPartialDisplay{S}{T}{V} \, \middle/ \, \myPartialDisplay{S}{T}{p} \right. \notag \\
				&= \frac{C_V / T}{C_p / T} \myTag{见\egref{EG_C_p-C_V}} \\
				&= \frac{C_V}{C_p} \fullstop
			\end{align}
			
			\begin{myProof}
				这里证明几个上面用到的结论：%TODO：20160402 偏微分关系的证明
				\begin{myEnum2}
					\item
						\begin{equation} \label{EQ_PARTIAL_DIFFERENTIAL_TO_JACOBI_DET}
							\myPartial{f}{g}{h} = \frac{ \pd \; (f, \, h) }{ \pd \; (g, \, h) } \fullstop
						\end{equation}
					\item
						\begin{equation} \label{EQ_JACOBI_DET_FRACTION}
							\myPartial{f}{g}{h} = \left. \frac{ \pd \; (f, \, h) }{ \pd \; (x, \, y) } \, \middle/ \, \frac{ \pd \; (g, \, h) }{ \pd \; (x, \, y) } \right. \fullstop
						\end{equation}
					\item
						\begin{equation} \label{EQ_INVERSE_JACOBI_DET}
							\frac{ \pd \; (f, \, g) }{ \pd \; (x, \, y) } = \left[ \frac{ \pd \; (x, \, y) }{ \pd \; (f, \, g) } \right]^{-1} \fullstop
						\end{equation}
				\end{myEnum2}
			\end{myProof}
		\end{myExample}
		
\section{\texorpdfstring{Joule--\!\! Thomson}{Joule--Thomson}效应；绝热膨胀与制冷}%HACK：20160625 em dash 的 kerning
	\subsection{节流过程}
		为了研究气体的内能，Joule先采用了气体自由膨胀的方法［见\secref{SEC_热容与焓；理想气体的性质}\subsecref{SUBSEC_理想气体的性质}］，但其精度不佳。之后（1852年），他又与Thomson采取了另外的手段，即\emphA{节流过程}。
		
		如图，%TODO：20160402 图片，节流过程
		在一根绝热管中间放置一个多孔塞，使气体无法很快地通过。保持多孔塞两边的压强差，则气体将从其一边不断地流到另一边，该过程就是\emphB{节流过程}。
		
		开始时，气体完全在多孔塞的一边。设其压强、体积、内能分别为 $p_1$、$V_1$ 和 $U_1$。让外界做功，直到气体完全流到多孔塞的另一边。此时，气体的压强、体积、内能分别为 $p_2$、$V_2$ 和 $U_2$。显然，外界做功
		\begin{equation}
			\db W = p_1 V_1 - p_2 V_2 \semicomma
		\end{equation}
		因为是绝热过程，因此 $\db Q = 0$。根据热力学第一定律，
		\begin{equation}
			U_2 - U_1 = \db W + \db Q = p_1 V_1 - p_2 V_2 \comma
		\end{equation}
		即
		\begin{equation}
			U_1 + p_1 V_1 = U_2 + p_2 V_2 \fullstop
		\end{equation}
		这说明节流过程是一个\emphB{等焓过程}，即 $H_1 = H_2$。\footnote{
			要注意该过程是一个\emphB{不可逆过程}。
		}
		
		定义\emphA{Joule--Thomson系数}
		\begin{equation}
			\m \eqdef \myPartial{T}{p}{H} \fullstop
		\end{equation}
		根据\emphB{偏导数三乘积法则} \eqref{EQ_CYCLIC_RELATION} 式，可知
		\begin{equation}
			\m = -\myPartial{H}{p}{T} \myPartial{T}{H}{p} \fullstop
			%TODO：20160402 没写，制冷曲线等
		\end{equation}
		
	\subsection{绝热膨胀过程}
		对于绝热过程，选取 $T$ 和 $p$ 作为独立变量，则
		\begin{equation}
			\dd S = \myPartial{S}{T}{p} \dd T+ \myPartial{S}{p}{T} \dd p = 0 \fullstop
		\end{equation}
		因此
		\begin{align}
			\myPartial{T}{p}{S} &= -\myPartial{S}{p}{T} \myPartial{T}{S}{p} \myTag{偏导数三乘积法则} \\
			&= -\left. \myPartialDisplay{S}{p}{T} \, \middle/ \, \myPartialDisplay{S}{T}{p} \right. \myTag{倒数关系} \\
			&= -\frac{T}{C_p} \myPartial{S}{p}{T} \myTag{见式~\eqref{EQ_C_p_IN_T_AND_S}} \\
			&= \frac{T}{C_p} \myPartial{V}{T}{p} \myTag{Maxwell关系} \\
			&= \frac{T V \a}{C_p} > 0  \myTagNumbering{见式~\eqref{EQ_DEFINITION_OF_EXPANSION_COEFFICIENT}} \fullstop
		\end{align}
		随着压强的增大，分子间距离减小，相互作用力增大，因而总能量减少，分子平均动能增加，从而使得温度上升。%TODO：20160402 解释对不对——能量
	\subsection{低温的实现}
	
\section{热力学函数的确定}
\section{特性函数}
\section{热辐射的热力学理论}
	这里所讲的热辐射即\emphB{黑体辐射}。温度升高，电偶极子发生振荡，因而辐射出电磁场（电磁波）。这就是热辐射的来源。
	
	本节内容说明热力学不仅适用于实物粒子，还适用于\emphB{场}。
	
	\subsection{内能密度}
		定义\emphA{内能密度}
		\begin{equation}
			u \eqdef \frac{U}{V}
		\end{equation}
		代表热辐射单位体积的内能。下面证明它是温度 $T$ 的\emphA{普适函数}。
		
		设有A、B两个空腔，其中充满了热辐射。A、B具有相同的温度 $T$。它们之间由一个细管连接，细管中有一个滤波片，它只允许频率在 $\n$ 到 $\n + \dd \n$ 之间的辐射通过。
		
		%TODO：20160406 内能密度证明，图片
		
		因此
		\begin{equation}
			u_\text{A}(\n) = u_\text{B}(\n)
		\end{equation}
		于是
		\begin{equation}
			u_\text{A} = u_\text{B} \fullstop
		\end{equation}
		该结果与A、B两空腔的形状、大小、腔壁材质均无关，这就说明了 $u$ 是 $T$ 的普适函数。因此，热辐射的内能
		\begin{equation} \label{EQ_RADIATION_INTERNAL_ENERGY}
			U = U(T, \, V) = V u(T) \fullstop
		\end{equation}
		
	\subsection{压强公式}
		根据实验结果，辐射压强
		\begin{equation} \label{EQ_RADIATION_PRESSURE}
			p = \frac{1}{3} u \fullstop
		\end{equation}
		该结果也可以通过电磁理论证明得到。
		%TODO：20160406 辐射压强的证明
		
	\subsection{热力学函数}
		与一般的 $p\text{-}V\text{-}T$ 系统一样，对于热辐射，有
		\begin{equation}
			\dd U = T \dd Ｓ - p \dd V \fullstop
		\end{equation}
		因此
		\begin{align}
			\myPartial{U}{V}{T} &= T \myPartial{S}{V}{T} - p \notag \\
			&= T \myPartial{p}{T}{V} - p \myTagNumbering{Maxwell关系} \fullstop
		\end{align}
		代入式~\eqref{EQ_RADIATION_INTERNAL_ENERGY} 和式~\eqref{EQ_RADIATION_PRESSURE}，可得
		\begin{equation}
			u = T \left( \frac{1}{3} \frac{\dd u}{\dd T} \right) - \frac{u}{3} \comma
		\end{equation}
		即
		\begin{equation}
			\frac{\dd u}{u} = \frac{4 \dd T}{T}　\fullstop
		\end{equation}
		积分可得
		\begin{equation}
			u = a T^4 \comma
		\end{equation}
		其中的 $a$ 是积分常数。于是热辐射的内能
		\begin{equation}
			U = a T^4 V \fullstop
		\end{equation}
		显然，$V =0 $ 时也成立 $U = 0$，因此上式不含附加的常数。同时，还可以得到压强与温度的关系：
		\begin{equation} \label{EQ_p_T_RELATION_OF_RADIATION}
			p = \frac{1}{3} u = \frac{1}{3} a T^4 \fullstop
		\end{equation}
		
		对于熵，
		\begin{align}
			\dd S &= \frac{1}{T} (\dd U + p \dd V) \notag \\
			&= \frac{1}{T} \left[ \dd \; (u V) + \frac{1}{3} u \dd V \right] \notag \\
			&= \frac{1}{T} \dd \; (a T^4 V^2) + \frac{1}{3 T} a T^4 V \dd V \notag \\
%			&= \frac{a}{T} (4 T^3 V^2 \dd T + 2 T^4 V \dd V) + \frac{1}{3} a T^3 V \ddV \notag \\
			&= \dd \left( \frac{4}{3} a T^3 V \right) \fullstop
		\end{align}
		积分得
		\begin{equation}
			S = \frac{4}{3} a T^3 V + S_0 \comma
		\end{equation}
		这里的 $S_0$ 是积分常数。当 $V \approach 0$ 时，显然应该有 $S \approach 0$。因此$S_0 = 0$，即
		\begin{equation} \label{EQ_RADIATION_ENTROPY}
			S = \frac{4}{3} a T^3 V \fullstop
		\end{equation}
		
		利用内能和熵的表达式，可以求出其他几个热力学函数：
		\begin{braceEq}
			H &= U + p V = \frac{4}{3} a T^4 V \semicomma \\
			F &= U - T S = -\frac{1}{3} a T^4 V \semicomma \\
			G &= U - T S + p V = 0 \fullstop
		\end{braceEq}
		这里的 $G = 0$ 有重要意义：热辐射的化学势为零，微观上说明光子数不守恒。
		
		根据式~\eqref{EQ_RADIATION_ENTROPY}，可知
		\begin{braceEq}
			C_V &= T \myPartial{S}{T}{V} = T \cdot 4 a T^2 V = 3 S \semicomma \\
			C_p &= T \myPartial{S}{T}{p} = \infty \fullstop %TODO：20160413 为什么C_p=无穷
		\end{braceEq}
		$C_p = \infty$，说明若 $p$ 不变，则 $T$ 也将保持不变。
		
		考虑可逆绝热过程，即 $S = \const$。如果选取 $T$ 和 $V$ 作为变量，根据式~\eqref{EQ_RADIATION_ENTROPY}，则有
		\begin{equation}
			T^3 V = \const \semicomma
		\end{equation}
		若选取 $p$ 和 $V$ 作为变量，利用式~\eqref{EQ_p_T_RELATION_OF_RADIATION}，便可得到
		\begin{equation} \label{EQ_STATE_EQUATION_OF_ADIABATIC_PROCESS_OF_RADIATION}
			p V^{4/3} = \const
		\end{equation}
		该式与理想气体可逆绝热过程方程 \eqref{EQ_STATE_EQUATION_OF_ADIABATIC_PROCESS_IN_p_V} 式（$p V^\g = \const$）形式上很类似，故它也被称作“光子气体”。但对于理想气体，$\g = C_p / C_V$，而式~\eqref{EQ_STATE_EQUATION_OF_ADIABATIC_PROCESS_OF_RADIATION}中的 $4/3$ 只是幂指数，与 $\g$ 无关。对于热辐射，$\g$ 实际上等于 $\infty$。
		
		\raggedbottom%FIXME：20160315 交叉引用、脚注每页重新计数失效，必须加上该行
		\pagebreak
			
		\chapter{相变的热力学理论}
			我们所研究的系统是逐渐复杂的。首先是\emphB{独立子体系}，如理想气体、Fermi气体、Bose气体等；之后是\emphB{近独立子体系}，如准粒子气体；%TODO：20160413 翻译
本章主要探讨\emphB{相互作用体系}，它包含多种物态，也会有相变的发生。

\section{热动平衡判据；开系热力学}
	在 \secref{SEC_熵增加原理与最大功}中我们已经证明，对于孤立体系（即绝热过程），熵将不断增加，直至达到极大（平衡）。因此，熵达到极大便可以作为体系达到平衡的判据。
	
	\subsection{平衡判据}
		\begin{myEnum1}
			\myItem{熵判据}
				由上，在 $\vd U = 0, \, \vd V = 0, \, \vd N = 0$ 的前提下，$\vd S = 0, \, \vd^{\:2} \! S <0$ 即说明（孤立）体系达到了平衡。这里的“$\vd\,$”表示\emphB{虚变化}，与“$\dd\,$”代表的\emphB{真实变化}有所不同，它与虚功原理中的虚位移是类似的。%HACK：20160420 δS的缩进
				
				平衡分为三种：稳定平衡、亚稳平衡和不稳定平衡。如前所述，熵的极大值对应平衡态。稳定平衡对应其中\emphB{最大}的极大值，而亚稳平衡对应其他较小的极大值，即对于无限小的变动是稳定的，而对于有限的变动则是不稳定的。不稳定平衡对应极小值，虽然有 $\vd S = 0$，但却不满足 $\vd^{\:2} \! S <0$。平衡的稳定性可以用力学类比来理解。重力势能 $E_\text{p}$ 的极小值对应平衡。其最小值对应稳定平衡；但相对极小对于大的扰动是不稳定的，所以对应亚稳平衡。而极大值则对应不稳定平衡，稍有扰动就会偏离。%TODO：20160420 图片
				
				对于 $\vd S = 0$、$\vd^{\:2} \! S =0$ 的临界态，将 $S$ 围绕极值点Taylor作展开，可得
				\begin{equation}
					\tl{\incr\,} S = \vd S + \frac{1}{2!} \vd^{\:2} \! S + \frac{1}{3!} \vd^{\:3} \! S + \frac{1}{4!} \vd^{\:4} \! S + \cdots < 0 \fullstop
				\end{equation}
				式中的“$\tl{\incr\,}$”同样表示虚变动。由于要满足 $S \rightarrow -S$ 的对称性，因此 $\vd^{\:3} \! S = 0$。从而%TODO：20160420 三阶项为何等于零
				\begin{equation}
					\vd^{\:4} \! S < 0 \comma
				\end{equation}
				这就是临界态的平衡判据。
				
			\myItem{自由能判据}
				考虑一个系统与热库（即环境）组成的复合体系，它是一个孤立系。显然，总内能
				\begin{equation}
					U_0 = U_\text{sys} + U_\text{res} = \const \comma
				\end{equation}
				总体积
				\begin{equation}
					V_0 = V_\text{sys} + V_\text{res} = \const
				\end{equation}
				设想体系发生了一个虚变动，则
				\begin{braceEq}[gather]
					\vd U_\text{sys} + \vd U_\text{res} = 0 \comma \\
					\vd V_\text{sys} + \vd V_\text{res} = 0 \fullstop \label{EQ_DELTA_V_SYS+DELTA_V_RES=0}
				\end{braceEq}
				根据熵判据，
				\begin{braceEq}
					\vd S_0 &= \vd \; (S_\text{sys} + S_\text{res}) = 0 \comma \\
					\vd^{\:2} \! S_0 &= \vd^{\:2} (S_\text{sys} + S_\text{res}) < 0 \fullstop
				\end{braceEq}
				根据热力学基本微分方程 \eqref{EQ_FUNDAMENTAL_EQUATION_FOR_PVT_SYSTEM} 式，可得
				\begin{equation}
					\vd U_\text{res} = T_\text{res} \vd S_\text{res} + p_\text{res}  \vd V_\text{res} \fullstop \footnote{
						已经把式~\eqref{EQ_FUNDAMENTAL_EQUATION_FOR_PVT_SYSTEM} 中的微分“$\dd\,$”改成了变分“$\vd\,$”的形式。
					}
				\end{equation}
				
				在系统温度、体积均不变（即 $\vd T_\text{sys} = 0$、$\vd V_\text{sys} = 0$）的情形下，根据式~\eqref{EQ_DELTA_V_SYS+DELTA_V_RES=0} 可知 $\vd V_\text{res} = 0$。由于是平衡态，又有 $T_\text{res} = T_\text{sys}$。因此
				
				\begin{align}
					\vd F_\text{sys} &= \vd \; (U_\text{sys} - T_\text{sys} S_\text{sys}) \\
					&= \vd U_\text{sys} - T_\text{sys} \vd S_\text{sys} \\
					&= -\vd U_\text{res} + T_\text{res} \vd S_\text{res} \\
					&= -p_\text{res}  \vd V_\text{res} \\
					&= -p_\text{res}  \vd V_\text{sys} = 0 \comma \\
					\vd^{\:2} \! F_\text{sys} &= %TODO：20160420 推导过程
				\end{align}
				
			\myItem{Gibbs函数判据}
			\myItem{内能判据}
		\end{myEnum1}
		
	\subsection{开系热力学}
		所谓“开系”，即\emphA{开放系（open system）}，它指粒子数可变且有能量交换的系统。
		
		对于封闭系，我们已经知道
		\begin{equation}
			\dd G = -S \dd T + V \dd p \fullstop
		\end{equation}
		而对于开放系，需引进\emphA{化学势} $\m$，其定义为
		\begin{equation}
			\m \eqdef \myPartial{G}{n}{T, \, p} \fullstop \footnote{这里的$n$指物质的量，而非粒子数。定义$\m \eqdef (\pd G / \pd N)_{T,\,p}$ 当然也可以，不过我们只采用第一种定义。}
		\end{equation}
		因此，
		\begin{equation}
			\dd G = -S \dd T + V \dd p + \m \dd n \fullstop
		\end{equation}
		利用Legendre变换，可得粒子数可变系统的热力学基本方程：
		\begin{braceEq}
			\dd U &= T \dd S - p \dd V + \m \dd n \comma \\
			\dd H &= T \dd S + V \dd p + \m \dd n \comma \\
			\dd F &= -S \dd T - p \dd V + \m \dd n \fullstop
		\end{braceEq}
		于是可以得出 $\m$ 的几个等价定义：
		\begin{equation}
			\m = \myPartial{U}{n}{S, \, V} = \myPartial{H}{n}{S, \, p} = \myPartial{F}{n}{T, \, V} = \myPartial{G}{n}{T, \, p} \fullstop
		\end{equation}
		
		Gibbs函数 $G(T,\,p,\,n)$ 是广延量。因此可定义\emphA{摩尔Gibbs函数} $G_\text{m}$，使得
		\begin{equation}
			G(T,\,p,\,n) = n G_\text{m} (T,\,p) \fullstop
		\end{equation}
		因此，
		\begin{equation}
			\m \eqdef \myPartial{G}{n}{T, \, p} = \left[ \frac{\pd}{\pd n} (n G_\text{m}) \right]_{T, \, p} = G_\text{m} \fullstop
		\end{equation}
		对于 \SI{1}{\mol} 物质，有
		\begin{equation} \label{EQ_FUNDAMENTAL_EQUATION_WITH_CHEMICAL_POTENTIAL}
			\dd \m = \dd G_\text{m} = -S_\text{m} \dd T + V_\text{m} \dd p \comma
		\end{equation}
		这是关于 $\m$ 的基本微分方程。
		
		\blankline
		
		定义\emphA{巨势（grand potential）}或\emphA{巨热力学势}
		\begin{equation} \label{EQ_DEF_OF_GRAND_POTENTIAL}
			\Y \eqdef F -\m n = F - G = -p \dd V \comma
		\end{equation}
		则其微分
		\begin{align}
			\dd \Y &= \dd F - \dd \; (\m n) \notag \\
			&= -S \dd T - p \dd V + \m \dd n - (\m \dd n + n \dd \mu) \notag \\
			&= -S \dd T - p \dd V - n \dd \mu \fullstop \footnotemark
		\end{align} \footnotetext{根据式~\eqref{EQ_DEF_OF_GRAND_POTENTIAL}，可知 $\dd \Y = -\dd \; (p V)$。实际上，
		\begin{align*}
			-S \dd T - p \dd V - n \dd \mu &= -S \dd T - p \dd V - n (-S_\text{m} \dd T + V_\text{m} \dd p) \\
			&= -S \dd T - p \dd V + S \dd T - V \dd p = -\dd \; (p V) \comma
		\end{align*}
		这是不矛盾的。}
		在统计物理中，巨势与巨配分函数有关。
		
\section{平衡条件与稳定条件} \label{SEC_平衡条件与稳定条件}
	\subsection{平衡条件} \label{SUBSEC_平衡条件}
		对于一个\emphB{单元两相}（即一种组分，两种状态）的系统，其平衡条件为
		\begin{equation}
			\vd S = \vd S_1 + \vd S_2 = 0 \comma
		\end{equation}
		其中的下标“1”和“2”分别表示两个相。根据热力学基本方程，有
		\begin{braceEq}
			\vd S_1 &= \frac{1}{T_1} (\vd U_1 + p_1 \vd V_1 - \m_1 \vd n_1) \comma \\
			\vd S_2 &= \frac{1}{T_2} (\vd U_2 + p_2 \vd V_2 - \m_2 \vd n_2) \fullstop
		\end{braceEq}
		设整个系统是孤立的，则
		\begin{braceEq}
			\vd U_1 + \vd U_2 &= 0 \comma \\
			\vd V_1 + \vd V_2 &= 0 \comma \\
			\vd n_1 + \vd n_2 &= 0 \fullstop
		\end{braceEq}
		因此
		\begin{align}
			\vd S &= \vd S_1 + \vd S_2 \notag \\
			&= \left[ \frac{1}{T_1} (\vd U_1 + p_1 \vd V_1 - \m_1 \vd n_1) \right] + \left[ \frac{1}{T_2} (\vd U_2 + p_2 \vd V_2 - \m_2 \vd n_2) \right] \notag \\
			&= \left[ \frac{1}{T_1} (\vd U_1 + p_1 \vd V_1 - \m_1 \vd n_1) \right] + \left[ \frac{1}{T_2} (-\vd U_2 - p_2 \vd V_2 + \m_2 \vd n_2) \right] \notag \\
			&= \vd U_1 \left( \frac{1}{T_1} - \frac{1}{T_2} \right) + \vd V_1 \left( \frac{p_1}{T_1} - \frac{p_2}{T_2} \right) - \vd n_1 \left( \frac{\m_1}{T_1} - \frac{\m_2}{T_2} \right) = 0 \fullstop
		\end{align}
		由于 $U_1$、$V_1$ 和 $n_1$ 是相互独立的，所以它们的系数都应该等于零，即
		\begin{braceEq}
			T_1 &= T_2 \comma \quad \text{（热平衡条件）} \\
			p_1 &= p_2 \comma \quad \text{（力学平衡条件）} \\
			\m_1 &= \m_2 \fullstop \quad \text{（化学平衡条件）}%TODO：20160420 相变平衡条件
		\end{braceEq}
		若 $T$ 不等，则仍有能量流动；$p$ 不等，则两相界面可以移动；$\m$ 不等，则每一相中的粒子数仍在变化，因此都不是平衡。
		
	\subsection{稳定条件}
		计算熵的二级变分：
		\begin{align}
			\vd^{\:2} \! S_1 = \frac{1}{2} \left( \frac{\pd^{\:2} \! S_1}{\pd U_1^2} \right) %TODO：20160420 (δS)^2还是δ^2 S
		\end{align}
		
\section{单元系的复相平衡}%；相变分类；Ehrenfest公式}
	\subsection{单元系的相图}
		设单元系中有两个相（$\a$ 和 $\b$），根据\secref{SEC_平衡条件与稳定条件}\subsecref{SUBSEC_平衡条件}，其平衡条件为
		\begin{braceEq}
			T^\a &= T^\b \comma \\
			p^\a &= p^\b \comma \\
			\m^\a &= \m^\b \fullstop
		\end{braceEq}
		令 $T$、$p$ 分别为两相共同的温度和压强，并取它们为独立变量，则相变的平衡条件为
		\begin{equation} \label{EQ_PHASE_TRANSFORMATION_CONDITION_FOR_1_COMPOSITION_AND_2_PHASE}
			\m^\a (T, \, p) = \m^\b (T, \, p) \fullstop
		\end{equation}
		这实际上就是 $p$-$T$ 平面内一条曲线的方程。
		
		如果存在三个相（$\a$、$\b$ 和 $\g$）并且同时达到平衡，则有
		\begin{equation}
			\m^\a (T, \, p) = \m^\b (T, \, p) = \m^\g (T, \,p) \fullstop
		\end{equation}
		这表示 $p$-$T$ 平面内的一个点，即\emphA{三相点}。
		
		\begin{figure}[h]
			\centering
%			\includegraphics[width = 6 cm]{Phase_diagram_of_water.svg}
			\caption{水的相图}
			\label{FIG_PHASE_DIAGRAM_OF_WATER}
		\end{figure}
		
		单元系的平衡性质利用\emphA{相图}可以清楚地表示。图%~\ref{FIG_PHASE_DIAGRAM_OF_WATER} 
		显示的是水的相图。在三相点处，水的固、液、气三相共存。该点对应的温度为 \SI{273.16}{\kelvin}，压强为 \SI{611.657}{\pascal}。另一个值得注意的点是\emphA{临界点}，气液两相的平衡曲线终结于此处。对于水，该点对应的温度为 \SI{647}{\kelvin}，压强为 \SI{22.064}{\mega\pascal}。
		%TODO：20160711 相图
		
	\subsection{Clausius--Clapeyron方程}
		由式~\eqref{EQ_PHASE_TRANSFORMATION_CONDITION_FOR_1_COMPOSITION_AND_2_PHASE}，单元系中两相平衡的条件为
		\begin{equation}
			\m^\a (T, \, p) = \m^\b (T, \, p) \fullstop
		\end{equation}
		设 $(T + \dd T, \, p + \dd p)$ 为平衡曲线上邻近 $(T, \, p)$ 的一点。当 $T$ 和 $p$ 沿平衡曲线变化到 $(T + \dd T, \, p + \dd p)$ 时，$\m^\a$ 和 $\m^\b$ 仍将满足平衡条件，即
		\begin{equation}
			\m^\a (T + \dd T, \, p + \dd p) = \m^\b (T + \dd T, \, p + \dd p) \fullstop
		\end{equation}
		两边都在点 $(T, \, p)$ 处作Taylor展开，可得
		\begin{equation}
			\m^\a (T, \, p) + \dd \m^\a = \m^\b (T, \, p) + \dd \m^\b \fullstop
		\end{equation}
		因此有
		\begin{equation} \label{EQ_dμα=dμβ}
			\dd \m^\a = \dd \m^\b \fullstop
		\end{equation}
		根据关于 $\m$ 的基本微分方程 \eqref{EQ_FUNDAMENTAL_EQUATION_WITH_CHEMICAL_POTENTIAL} 式，可知
		\begin{braceEq}
			\dd \m^\a &= -S^\a_\text{m} \dd T + V^\a_\text{m} \dd p \comma \\
			\dd \m^\b &= -S^\b_\text{m} \dd T + V^\b_\text{m} \dd p \fullstop
		\end{braceEq}
		将其代入式~\eqref{EQ_dμα=dμβ}，便有
		\begin{equation} \label{EQ_CLAUSIUS_CLAPEYRON_RELATION_1}
			\frac{\dd p}{\dd T} = \frac{\displaystyle S^\a_\text{m} - S^\b_\text{m}}{\displaystyle V^\a_\text{m} - V^\b_\text{m}} \fullstop
		\end{equation}
		
		定义\emphA{相变潜热}
		\begin{equation}
			L_{\a \b} = T \left( S^\a_\text{m} - S^\b_\text{m} \right) \comma
		\end{equation}
		它表示 \SI{1}{\mol} 物质在保持温度、压强不变的条件下由 $\a$ 相转变为 $\b$ 相所吸收的热量。这样，\eqref{EQ_CLAUSIUS_CLAPEYRON_RELATION_1} 就可以表述为
		\begin{equation}
			\frac{\dd p}{\dd T} = \frac{L_{\a \b}}{\displaystyle T \left( V^\a_\text{m} - V^\b_\text{m} \right)} \comma
		\end{equation}
		此即\emphA{Clausius--Clapeyron方程}。
		
		通常情况下，液体到气体的相变满足 $\dd p / \dd T > 0$；而液体到固体的相变满足 $\dd p / \dd T < 0$。需注意，水是一个常见的例外。液态的水变成冰，$\dd p / \dd T > 0$。
		
	\subsection{相变的分类}
		根据Ehrenfest的分类方案，$n$ 级相变对应Gibbs函数 $n$ 阶导数的不连续变化。相变点处，化学势连续，而化学势一阶导数不连续的相变称为\emphA{一级相变}；化学势及其一阶导数均连续，而二阶导数不连续的相变称为\emphA{二级相变}。
		
		Fisher推广了Ehrenfest的分类方案。在他的方案中，化学势及其一阶导数均连续，而二阶导数或者不连续、或者\emphB{发散}的相变称为二级相变。
		
		……%TODO：20160711 举例
		
	\subsection{Ehrenfest方程}
		Clausius--Clapeyron方程的一种形式是 \eqref{EQ_CLAUSIUS_CLAPEYRON_RELATION_1} 式：
		\begin{equation} \label{EQ_CLAUSIUS_CLAPEYRON_RELATION_BEFORE_EHRENFEST_EQUATION}
			\frac{\dd p}{\dd T} = \frac{\displaystyle S^\b_\text{m} - S^\g_\text{m}}{\displaystyle V^\b_\text{m} - V^\g_\text{m}} = \frac{\incr S_\text{m}}{\incr V_\text{m}}\fullstop
		\end{equation}
		对于二级相变（按照Ehrenfest的分类方案），$\incr S_\text{m} = 0$ 且 $\incr V_\text{m} = 0$，因而上式成为 $0/0$ 的不定式。这时，按照l'Hôpital法则处理，即将分子分母同时对 $T$ 求偏导，可得
		\begin{equation} \label{EQ_EHRENFEST_EQUATION_A}
			\frac{\dd p}{\dd T}
			=\frac{\incr \myPartialDisplay{\displaystyle S_\text{m}}{T}{p} }{\incr \myPartialDisplay{\displaystyle V_\text{m}}{T}{p} }
			=\frac{\incr \left( C_{\text{m},\,p} / T \right) }{\incr \left(V_\text{m} \a\right) }
			=\frac{\displaystyle \incr C_{\text{m},\,p}}{T V_\text{m} \incr\a} \comma
		\end{equation}
		其中的 $\a$ 为膨胀系数，而 $V_\text{m}=V^\b_\text{m}=V^\g_\text{m}$。
		
		若是将式~\eqref{EQ_CLAUSIUS_CLAPEYRON_RELATION_BEFORE_EHRENFEST_EQUATION} 的分子分母同时对 $p$ 求偏导，则将得到
		\begin{equation} \label{EQ_EHRENFEST_EQUATION_B}
			\frac{\dd p}{\dd T}
			=\frac{\incr \myPartialDisplay{\displaystyle S_\text{m}}{p}{T} }{\incr \myPartialDisplay{\displaystyle V_\text{m}}{p}{T} }
			=\frac{-\incr \myPartialDisplay{\displaystyle V_\text{m}}{T}{p} }{\incr \myPartialDisplay{\displaystyle V_\text{m}}{p}{T} }
			=\frac{-\incr \left(V_\text{m} \a \right) }{\incr \left(-V_\text{m} \k_T \right) }
			=\frac{\a}{\incr\k_T} \comma
		\end{equation}
		其中的 $\k_T$ 为等温压缩系数。另外，这里的推导用到了Maxwell关系\eqref{EQ_MAXWELL_RELATION_FROM_G} 式。
		
		式~\eqref{EQ_EHRENFEST_EQUATION_A} 和 \eqref{EQ_EHRENFEST_EQUATION_B} 都称为\emphA{Ehrenfest方程}。要使其自洽，显然应该有
		\begin{equation}
			\frac{\displaystyle \incr C_{\text{m},\,p}}{T V_\text{m} \incr\a}
			=\frac{\a}{\incr\k_T} \comma
		\end{equation}
		即
		\begin{equation}
			\incr C_{\text{m},\,p} = T V_\text{m} \frac{\left(\incr\a\right)^2}{\incr\k_T} \fullstop
		\end{equation}
\section{气液相变；临界点行为}
\section{临界指数；普适性}
\section{Landau平均场理论}
	
\raggedbottom%FIXME：20260325 交叉引用、脚注每页重新计数失效，必须加上该行
\pagebreak
			
%		\chapter{多元复相 热力学第三定律}
%			\section{多元系热力学函数及热力学方程}
\section{多元系的复相平衡}
\section{Gibbs相律}
\section{混合理想气体}
\section{热力学第三定律}

%			
%	\part{统计物理}
		\chapter{统计物理学基本概念}
			\section{微观态的经典及量子描述}
	\subsection{单粒子的经典描述}
		微观态的经典描述以经典力学为基础，通常采用广义坐标与广义动量的形式。
		
		对于一个有 $r$ 个自由度的系统，需要用 $2r$ 个变量来描述其运动状态，即 $r$ 个广义坐标和 $r$ 个广义动量：
		\begin{equation}
			(q_i, \, p_i) \quad (i = 1, \, 2, \, \cdots, r) \fullstop
		\end{equation}
		系统的Hamilton量为
		\begin{equation}
			H = H(q_1, \, q_2, \, \cdots, q_r; \, p_1, \, p_2, \, \cdots, p_r; \, t) \comma
		\end{equation}
		正则方程为
		\begin{braceEq}
			\dot{q_i} &= \frac{\pd H}{\pd p_i} \comma \\
			\dot{p_i} &= -\frac{\pd H}{\pd q_i} \fullstop
		\end{braceEq}
		
		坐标和动量 $(q_1, \, q_2, \, \cdots, q_r; \, p_1, \, p_2, \, \cdots, p_r)$ 张成了一个 $2r$ 维空间，称为\emphA{\itshape{μ}空间}，每一组坐标和动量描述的点称为\emphA{代表点}。
		
		\begin{myExample}[自由粒子]
			对于一个 $r = 3$ 的自由粒子，有
			\begin{braceEq}
				p_x = m \dot{x} \comma \\
				p_y = m \dot{y} \comma \\
				p_z = m \dot{z} \fullstop
			\end{braceEq}
			其 $\m$ 空间由 $(x, \, y, \, z; \, p_x, \, p_y, \, p_z)$ 张成。Hamilton量
			\begin{equation}
				H = \frac{1}{2m} \left( p_x^2 + p_y^2 +p_z^2 \right) \fullstop
			\end{equation}
			%TODO：20160504 示意图
		\end{myExample}
		
		\begin{myExample}[一维谐振子]
			质量为 $m$ 的物体受力 $F = -A x$ 的作用做简谐运动，其角频率
			\begin{equation}
				\o = \sqrt{\frac{A}{m}} \comma
			\end{equation}
			则其Hamilton量
			\begin{equation}
				H = \frac{p^2}{2m} + \frac{A}{2} x^2 = \frac{p^2}{2m} + \frac{1}{2} m \o^2 x^2 \fullstop
			\end{equation}
			
			若总能量一定，即 $H = E$，则
			\begin{equation}
				\frac{p^2}{2mE} + \frac{x^2}{2E / m \o^2} = 1 \comma
			\end{equation}
			这在 $\m$ 空间中表示一个椭圆（见图），其面积
			\begin{equation} \label{EQ_ELLIPSE_OF_HARMONIC_OSCILLATOR_IN_MU_SPACE}
				S_\text{椭圆} = \p \, a b = \p \cdot \sqrt{\frac{2E}{m \o^2}} \cdot \sqrt{2mE} = \frac{2 \p E}{\o} \fullstop
			\end{equation}
			
			\begin{figure}[ht]
				\centering
				
				\begin{asy}
					pair O = (0, 0), x_axes = (8, 0), y_axes = (0, 5);
					
					real a = 5.5, b = 3;
					
					pair p1 = (0.4, 1.8), p2 = (5, 3.8);
					pair p3 = (2, -0.4), p4 = (6, -3);
					
					draw(Label("$x$", EndPoint), (-x_axes)--x_axes, Arrow);
					draw(Label("$p$", EndPoint), (0, -4.2)--y_axes, Arrow);
					
					draw(ellipse(O, a, b), linewidth(1) + color1);
					
					draw(O--(a, 0), linewidth(1.2));
					draw(O--(0, b), linewidth(1.2));
					
					draw(Label("$\sqrt{2mE}$", EndPoint), p1--p2);
					draw(Label("$\sqrt{\dfrac{\displaystyle 2E}{\displaystyle m\omega^2}}$", EndPoint), p3--p4);
					
					label("$O$", O, SW);
				\end{asy}
				\caption{$\m$ 空间中的一维谐振子}
				\label{FIG_1D_HARMONIC_OSCILLATOR_IN_MU_SPACE}
			\end{figure}
		\end{myExample}
		
		\begin{myExample}[转子]
			如图~\ref{FIG_ROTATOR} 所示，质量为 $m$ 的物体被轻杆连接在 $O$ 点处，可绕 $O$ 点运动。其Hamilton量为
			\begin{equation} \label{EQ_ROTATOR_HAMILTONIAN_1}
				H = \frac{m}{2} \left( \dot{x}^2 + \dot{y}^2 + \dot{z}^2 \right) \fullstop
			\end{equation}
			取球坐标系，则
			\begin{braceEq}
				x &= r \sin \th \cos \comma \\
				y &= r \sin \th \sin \vf \comma \\
				z &= r \cos \th \fullstop
			\end{braceEq}
			求导，则有
			\begin{braceEq}
				\dot{x} &= \dot{r} \sin\th \cos\vf + r \dot{\th} \cos\th \cos\vf - r \dot{\vf} \sin\th \sin\vf \\
				\dot{y} &= \dot{r} \sin\th \sin\vf + r \dot{\th} \cos\th \sin\vf + r \dot{\vf} \sin\th \cos\vf \\
				\dot{z} &= \dot{r} \cos\th - r \dot{\th} \sin\th \fullstop
			\end{braceEq}
			因此
			\begin{align}
				\dot{x}^2 + \dot{y}^2 + \dot{z}^2
				&= \left[ \dot{r}^2\sin^2\th + \left(r\dot{\th}\right)^2\cos^2\th + \left(r\dot{\vf}\right)^2\sin^2\th \right] \myTag{$\dot{x}^2$、$\dot{y}^2$ 平方项} \\
				&+ \left[ 2r\dot{r}\dot{\th}\sin\th\cos\th\cos^2\vf - 2r\dot{r}\dot{\vf}\sin^2\th\sin\vf\cos\vf - 2r^2\dot{\th}\dot{\vf}\sin\th\cos\th\sin\vf\cos\vf \right] \myTag{$\dot{x}^2$ 交叉项} \\
				&+ \left[ 2r\dot{r}\dot{\th}\sin\th\cos\th\sin^2\vf + 2r\dot{r}\dot{\vf}\sin^2\th\sin\vf\cos\vf + 2r^2\dot{\th}\dot{\vf}\sin\th\cos\th\sin\vf\cos\vf \right] \myTag{$\dot{y}^2$ 交叉项} \\
				&+ \left[ \dot{r}^2\cos^2\th - 2r\dot{r}\dot{\th}\sin\th\cos\th + \left(r\dot{\th}\right)^2\sin^2\th \right] \myTag{$\dot{z}^2$} \\
				&= \dot{r}^2 + \left(r\dot{\th}\right)^2 + \left(r\dot{\vf}\right)^2\sin^2\th \fullstop
			\end{align}
			代入 \eqref{EQ_ROTATOR_HAMILTONIAN_1}~式，得
			\begin{align}
				H &= \frac{m}{2} \left( \dot{r}^2 + r^2\dot{\th}^2 + r^2\dot{\vf}^2\sin^2\th \fullstop \right)
			\end{align}
			由于物体已被轻杆连接在了 $O$ 点，因而 $r$ 不变，$\dot{r}=0$。
			
			\begin{figure}[ht]
				\begin{asy}
					import my3D;
					
					triple O = (0, 0, 0), x_axes = (3.5, 0, 0), y_axes = (0, 5, 0), z_axes = (0, 0, 5);
					triple point_m = (2.5, 5, 6), point_n = (point_m.x, point_m.y, 0);
					
					pair O2 = project(O), point_m2 = project(point_m), point_n2 = project(point_n);
					
					draw(Label("$x$", EndPoint), O2--project(x_axes), Arrow);
					draw(Label("$y$", EndPoint), O2--project(y_axes), Arrow);
					draw(Label("$z$", EndPoint), O2--project(z_axes), Arrow);
					
					draw(Label("$\bm{r}$", MidPoint, black), O2--point_m2, linewidth(1.5) + color2);
					draw(O2--point_n2--point_m2, dashed + color1);
					fill(circle(point_m2, 0.2), color1);
					
					real angle_r = 0.7;
					draw(Label("$\varphi$", MidPoint, Relative(E)), angleMark(O, x_axes, point_n, angle_r));
					draw(Label("$\theta$", MidPoint, Relative((-1,-0.4))), angleMark(O, z_axes, point_m, angle_r));
					
					label("$O$", O2, (-1.5, 0.5));
					label("$m$", point_m2, (2, 0.5));
				\end{asy}
				\caption{转子的示意图}
				\label{FIG_ROTATOR}
			\end{figure}

			引入\emphA{共轭动量}
			\begin{braceEq}
				p_\th &= m r^2 \dot{\th} \comma \\
				p_\vf &= m r^2 \dot{\vf} \sin^2 \th \comma
			\end{braceEq}
			则系统的Hamilton量可写为
			\begin{equation}
				H = \frac{1}{2I} \left( p_\th^2 + \frac{1}{\sin^2 \th} p_\vf^2 \right) \comma
			\end{equation}
			其中的 $I = m r^2$ 是物体对 $O$ 点的\emphA{转动惯量}。
			
			在本例中，$\mu$ 空间由广义坐标和广义动量 $(\th, \, \vf; \, p_\th, \, p_\vf)$ 张成，它是四维的。
		\end{myExample}
		
	\subsection{单粒子的量子描述}
		微观态的量子描述以量子力学为基础。粒子的动量 $\vecb{p}$、能量 $E$ 满足\emphA{de Broglie关系}：
		\begin{braceEq}
			\vecb{p} &= \hbar \vecb{k} \comma \\
			E &= \hbar \o \comma
		\end{braceEq}
		其中的 $\hbar$ 称为\emphA{（约化）Planck常数}，其值为
		\begin{equation}
			\hbar = \frac{h}{2\p} = \SI{1.0545718e-34}{\joule\second} \fullstop
		\end{equation}
		式中的 $h$ 也是Planck常数。
		
		de Broglie关系说明微观粒子具有\emphA{波粒二象性}。这就引出了另一个重要结果——\emphA{不确定关系}\footnote{
			更精确的表述为 $\incr p \incr q \geqslant \hbar / 2 $。
		}：
		\begin{equation}
			\incr p \incr q \gtrsim h \fullstop
		\end{equation}
		由此，粒子的动量和坐标不可能被同时精确测量，因而其运动也就无法用经典的轨道概念描述，必须改用波函数来描述。
		
		粒子波函数 $\Y$ 满足的方程即\emphA{Schrödinger方程}：
		\begin{equation}
			\ii \frac{\pd}{\pd t} \Y = \op{H} \Y \comma
		\end{equation}
		式中的 $\op{H}$ 是Hamilton算符。在定态情况（即将时间变量分离后），Schrödinger方程化为
		\begin{equation}
			\op{H} \Y = E \Y \fullstop
		\end{equation}
		
		\begin{myExample}[立方体容器内的自由粒子]
			设粒子在边长为 $L$ 的立方体容器内运动，则其量子态（即波函数）有平面波的形式：
			\begin{equation}
				\Y_{n_1, \, n_2, \, n_3} (\vecb{r}) \sim \ee^{\ii \vecb{p} \cdot \vecb{r} / \hbar} \comma
				\quad n_i = \pm 1, \, \pm 2, \, \pm 3, \, \cdots \fullstop
			\end{equation}
			求解Schrödinger方程，可以发现动量与能量的本征值都是量子化的：
			\begin{braceEq}
				\vecb{p} &= p_x \vecb{i} + p_y \vecb{j} + p_z \vecb{k}
				= \frac{2\p\hbar}{L} n_1 \vecb{i} + n_2 \vecb{j} + n_3 \vecb{k} \comma \\
				E &= \frac{1}{2m} \left( p_x^2+p_y^2+p_z^2 \right)
				= \frac{\displaystyle 2\p^2 \hbar^2}{\displaystyle mL^2} \left( n_1^2+n_2^2+n_3^2 \right) \fullstop
			\end{braceEq} %HACK：20160624 乘方位置调整
			量子化的能量也成为\emphA{能级}。对于能级
			\begin{equation}
				E = \frac{\displaystyle 2\p^2 \hbar^2}{\displaystyle mL^2} \comma
			\end{equation}
			它对应6种量子态：
			\begin{equation}
				\left(\pm 1, \, 0, \, 0 \right),\,
				\left(0, \, \pm 1, \, 0 \right),\,
				\left(0, \, 0, \, \pm 1 \right) \fullstop
			\end{equation}
			这种现象称为能级\emphA{简并}。同一能级对应量子态的数目称为\emphA{简并度}。显然，这里的简并度为6。而能量更高的一个能级
			\begin{equation}
				E = \frac{\displaystyle 2\p^2 \hbar^2}{\displaystyle mL^2} \times 3
				=\frac{\displaystyle 2\p^2 \hbar^2}{\displaystyle mL^2} \times (1+1+1)
			\end{equation}
			则对应 $2^3=8$ 个量子态，它的简并度为8。
		\end{myExample}
		
		\begin{myExample}[一维谐振子]
			频率为 $\nu$ 的谐振子，其能量为
			\begin{equation}
				E_n = \left( n + \frac{1}{2} \right) h \nu \comma \quad n = 0, \, 1, \, 2, \, \cdots \fullstop
			\end{equation}
			可见该系统的简并度 $g = 1$。
			
			根据式~\eqref{EQ_ELLIPSE_OF_HARMONIC_OSCILLATOR_IN_MU_SPACE}，$\mu$ 空间中的椭圆面积为
			\begin{equation}
				S_n = \frac{2\pi E_n}{\o} =\frac{E_n}{\nu} = \left( n + \frac{1}{2} \right) h\fullstop
			\end{equation}
			因此两个相轨道之间的面积为 $h$，它对应一个量子态。
		\end{myExample}
		
		\begin{myExample}[转子]
			blablabla
		\end{myExample}%TODO：20160624 未完成
		
	\subsection{多粒子体系的描述}
		对于 $N$ 个粒子组成的体系，设每个粒子的自由度为 $r$，则每个粒子可用 $2r$ 个变量 $(q_1, \, q_2, \, \cdots, q_r; \allowbreak \, p_1, \, p_2, \, \cdots, p_r)$ 来描述。此时，系统的总自由度 %FIXME：20160624 公式换行
		\begin{equation}
			f = Nr \comma
		\end{equation}
		因而系统的运动需要用 $2f$ 个变量来刻画。它们张成了一个 $2f$ 维空间，称为\emphA{\itshape{Γ}空间}，也叫\emphA{相空间}。
%\section{}
%\section{}
%\section{}

\raggedbottom%FIXME：20260325 交叉引用、脚注每页重新计数失效，必须加上该行
\pagebreak

%		
%		\chapter{三大统计}
%			\include{Chapters/Chapter6}
%		
%		\chapter{系综理论}
%			\include{Chapters/Chapter7}
%		
%		\chapter{相变的统计物理简介}
%			\include{Chapters/Chapter8}
		
\end{document}